% !TEX program = pdflatex
% Long-form manuscript for Foundations of Physics / arXiv.
% Compiles with pdflatex + bibtex. See README_SUBMISSION.md for details.

\documentclass[11pt,a4paper,twoside]{article}

% ---------- encoding / fonts ----------
\usepackage[T1]{fontenc}
\usepackage[utf8]{inputenc}
\usepackage{lmodern}
\usepackage{microtype}

% ---------- math ----------
\usepackage{amsmath,amssymb,amsthm,mathtools}
\usepackage{bm}

% ---------- layout / typography ----------
\usepackage[a4paper,margin=2.8cm]{geometry}
\usepackage{setspace}
\onehalfspacing
\usepackage{parskip}
\usepackage{booktabs}
\usepackage{array}
\usepackage{longtable}
\usepackage{enumitem}
\setlist[itemize]{leftmargin=1.4em,itemsep=2pt,topsep=2pt}
\setlist[enumerate]{leftmargin=1.6em,itemsep=2pt,topsep=2pt}

% ---------- code ----------
\usepackage{xcolor}
\definecolor{codebg}{HTML}{F6F8FA}
\definecolor{codekw}{HTML}{0033B3}
\definecolor{codecomm}{HTML}{8E908C}
\definecolor{codestr}{HTML}{718C00}
\usepackage{listings}
\lstdefinelanguage{lean4}{%
  morekeywords={theorem,lemma,def,noncomputable,structure,class,instance,abbrev,example,namespace,end,section,import,open,variable,fun,let,in,if,then,else,match,with,where,have,show,by,from,intro,exact,apply,rfl,simp,rw,constructor,refine,use,cases,induction,unfold,set,calc,this,Prop,Type,Real,Complex,Nat,Int,Fin,Finset,sorry,axiom},%
  sensitive=true,%
  morecomment=[l]{--},%
  morecomment=[s]{/-}{-/},%
  morestring=[b]",%
}
\lstset{%
  basicstyle=\ttfamily\footnotesize,%
  backgroundcolor=\color{codebg},%
  keywordstyle=\color{codekw}\bfseries,%
  commentstyle=\color{codecomm}\itshape,%
  stringstyle=\color{codestr},%
  breaklines=true,%
  breakatwhitespace=true,%
  columns=fullflexible,%
  frame=single,%
  framerule=0.4pt,%
  xleftmargin=0pt,%
  xrightmargin=0pt,%
  aboveskip=6pt,belowskip=6pt,%
  language=lean4,%
  literate=%
    {→}{{$\to$}}1 {⇒}{{$\Rightarrow$}}1 {↦}{{$\mapsto$}}1
    {≤}{{$\le$}}1 {≥}{{$\ge$}}1 {≠}{{$\ne$}}1 {≈}{{$\approx$}}1
    {∧}{{$\wedge$}}1 {∨}{{$\vee$}}1 {¬}{{$\lnot$}}1 {∀}{{$\forall$}}1 {∃}{{$\exists$}}1
    {α}{{$\alpha$}}1 {β}{{$\beta$}}1 {γ}{{$\gamma$}}1 {δ}{{$\delta$}}1
    {ε}{{$\varepsilon$}}1 {ζ}{{$\zeta$}}1 {η}{{$\eta$}}1 {θ}{{$\theta$}}1
    {ι}{{$\iota$}}1 {κ}{{$\kappa$}}1 {λ}{{$\lambda$}}1 {μ}{{$\mu$}}1
    {ν}{{$\nu$}}1 {ξ}{{$\xi$}}1 {π}{{$\pi$}}1 {ρ}{{$\rho$}}1
    {σ}{{$\sigma$}}1 {τ}{{$\tau$}}1 {υ}{{$\upsilon$}}1 {φ}{{$\varphi$}}1
    {χ}{{$\chi$}}1 {ψ}{{$\psi$}}1 {ω}{{$\omega$}}1
    {Α}{{$A$}}1 {Β}{{$B$}}1 {Γ}{{$\Gamma$}}1 {Δ}{{$\Delta$}}1
    {Λ}{{$\Lambda$}}1 {Ξ}{{$\Xi$}}1 {Π}{{$\Pi$}}1 {Σ}{{$\Sigma$}}1
    {Φ}{{$\Phi$}}1 {Ψ}{{$\Psi$}}1 {Ω}{{$\Omega$}}1
    {ℝ}{{$\mathbb{R}$}}1 {ℂ}{{$\mathbb{C}$}}1 {ℕ}{{$\mathbb{N}$}}1 {ℤ}{{$\mathbb{Z}$}}1
    {ℓ}{{$\ell$}}1 {ℏ}{{$\hbar$}}1 {∑}{{$\sum$}}1 {∫}{{$\int$}}1 {⟨}{{$\langle$}}1 {⟩}{{$\rangle$}}1
    {‖}{{$\|$}}1 {∈}{{$\in$}}1 {∉}{{$\notin$}}1 {⊗}{{$\otimes$}}1 {⊕}{{$\oplus$}}1
    {·}{{$\cdot$}}1 {√}{{$\sqrt{}\!$}}1 {⊗ₜ}{{$\otimes_{\!t}$}}2
}

% ---------- theorem environments ----------
\newtheorem{theorem}{Theorem}[section]
\newtheorem{proposition}[theorem]{Proposition}
\newtheorem{corollary}[theorem]{Corollary}
\newtheorem{lemma}[theorem]{Lemma}
\theoremstyle{definition}
\newtheorem{definition}[theorem]{Definition}
\newtheorem{remark}[theorem]{Remark}

% ---------- references ----------
\usepackage[hidelinks,colorlinks=false,bookmarks=true,pdfusetitle]{hyperref}
\usepackage{cleveref}

% ---------- macros ----------
\newcommand{\KL}{I_{\mathrm{KL}}}
\newcommand{\Rbb}{\mathbb{R}}
\newcommand{\Cbb}{\mathbb{C}}
\newcommand{\Nbb}{\mathbb{N}}
\newcommand{\Zbb}{\mathbb{Z}}
\newcommand{\lP}{\ell_P}
\newcommand{\tP}{t_P}
\newcommand{\EP}{E_P}
\newcommand{\mP}{m_P}
\newcommand{\kB}{k_B}
\newcommand{\hbarc}{\hbar}
\newcommand{\Lmap}{\mathsf{L}}
\newcommand{\coarseGrain}{\operatorname{coarseGrain}}
\newcommand{\coarseGrainWithPhase}{\operatorname{coarseGrainWithPhase}}
\newcommand{\HasZF}{\mathsf{HasZeroFunctional}}
\newcommand{\CMM}{\mathsf{CommutatorMatchesMean}}
\newcommand{\xhat}{\hat{x}}
\newcommand{\phat}{\hat{p}}
\newcommand{\Lfield}{\mathsf{LatticeComplexField}}
\newcommand{\SnapSeq}{\mathsf{SnapshotSequence}}
\newcommand{\DynSeq}{\mathsf{DynamicalSnapshotSequence}}
\newcommand{\code}[1]{\texttt{#1}}
\newcommand{\BellField}{\operatorname{bellField}}
\newcommand{\Prop}{\mathsf{Prop}}

% ---------- title block ----------
\title{%
  \Large\bfseries
  Rigorous Machine-Checked Derivation of Non-Relativistic Quantum Mechanics\\
  from Discrete-Gravity Healing Dynamics
}

\author{%
  Norbert Marchewka\thanks{Correspondence: \texttt{norbert.marchewka44@gmail.com}}\\
  \small\emph{with the Omega-Theory V2 formalisation team (see Acknowledgments)}
}

\date{\today}

% ---------- document ----------
\begin{document}
\maketitle

\begin{abstract}\noindent
We derive \emph{seven pillars of non-relativistic quantum mechanics}
--- Schr\"odinger dynamics, the Born rule, the non-relativistic limit
of relativistic dispersion, two-slit interference, the Heisenberg
uncertainty principle, the measurement / collapse postulate, and
Tsirelson-bound-attaining entanglement with CHSH-inequality violation
--- together with a \emph{relativistic Klein--Gordon residue bound} on
the \emph{same} discrete-gravity substrate, all as machine-checked
Lean~4 theorems from the Planck-scale healing dynamics of
Omega-Theory V2. The substrate is a local update rule on $\Zbb^{4}$
whose equilibrium condition already reproduces Einstein's field
equations modulo a controlled $O(\lP)$ remainder (published
elsewhere); we show that the same substrate, once equipped with a
complex-valued tick-to-tick update rule, coarse-grains to standard
non-relativistic wave mechanics (with enough resolution to derive all
seven canonical QM postulates that textbooks take as axioms) and also
satisfies a relativistic Klein--Gordon residue bound
$\|\mathrm{KG}[\psi]\| \le 16/\lP^{2} + m^{2}c^{2}/\hbar^{2}$ in which
the Planck-scale UV cutoff appears explicitly as a prediction of the
theory's discrete-spacetime ontology.
\par\vspace{4pt}
The central construction is a map
$\Lmap:\SnapSeq\to\Lfield$ sending a time-indexed family of discrete
metrics to a complex field $\psi(x,n)$ on $\mathsf{LatticePoint}\times\Nbb$.
The amplitude is fixed by the Kullback--Leibler information density as
$|\psi|^{2}=\exp(-\KL)$; a phase-carrying extension
\code{coarseGrainWithPhase} admits plane-wave phases sourced from the
lattice geometry; a two-body tensor-product extension
\code{tensorProduct} carries the Bell-state construction.
\par\vspace{4pt}
All eight theorems (T1--T8 bundling Bell+CHSH, plus T9 Klein--Gordon)
and a capstone \code{grand\_qm\_emergence} theorem that bundles them
into a single nine-field $\Prop$-record
\code{QuantumMechanicsPostulates}, build clean in Lean~4 against
Mathlib v4.29.0 with zero \code{sorry} and no new axioms beyond the
eight physical Planck constants and the
Hildebrandt--Polthier--Wardetzky Laplacian--Ricci correspondence
(used only on the Einstein side, \emph{not} on the QM/KG bridge;
provably eliminable on three regimes). To the authors' knowledge,
this is the first machine-checked derivation in a proof assistant of
the full canonical-QM postulate set \emph{plus} the relativistic
Klein--Gordon residue bound from a \emph{single} discrete-gravity
substrate. Results that the substrate-derivation literature cites as
heuristics or axioms, we furnish as theorems.
\end{abstract}

\noindent\textbf{Status.} Draft, paper-in-preparation. Scope expanded
from 1-theorem to 7-theorem chain on 2026-04-15 (Option B plan
\code{PLAN\_QM\_BRIDGE.md} plus Phase 6A/B/C), then extended to the
nine-theorem relativistic bridge with the landing of Theorem~9
(\code{KleinGordon.lean}) and the Dirac stub
(\code{DiracOptional.lean}). All nine theorems landed as of
2026-04-15.\\
\textbf{Target venues.} \emph{Physical Review Letters} (4-page letter) or
\emph{Nature Physics} (6-page letter), with a long-form companion in
\emph{Foundations of Physics}. \emph{Journal of Mathematical Physics}
remains admissible if the letter-length venue is not reached.\\
\textbf{Formalisation.} Lean~4 + Mathlib v4.29.0. Source tree:
\path{PhysicsPapers/LeanFormalizationV2/}.

\tableofcontents

\section{Introduction}\label{sec:intro}

A recurring theme across discrete-substrate approaches to quantum
mechanics---Wolfram physics~\cite{Wolfram2020}, 't~Hooft's
cellular-automaton interpretation~\cite{tHooft2014,tHooft2016}, causal
sets~\cite{BombelliLeeMeyerSorkin1987}, stochastic
electrodynamics~\cite{dePenaCetto1996}---is that the smooth Schr\"odinger
evolution ought to emerge as the large-scale phenomenology of a discrete
local update rule. The difficulty is notorious: going from a local
combinatorial dynamics to a continuum wave equation has historically
required either informal coarse-graining arguments or postulates that
sneak in the Hilbert-space structure they set out to derive.

This paper reports a \emph{different kind} of contribution. We do not
propose a new physical mechanism. We take the mechanism that is already
built into the Omega-Theory V2 Lean formalization---a Planck-scale
lattice with a healing flow that stabilises the metric against
computational truncation errors, and from which Einstein's equations
have already been derived as a theorem at equilibrium---and we show, as
a sequence of Lean~4 theorems, that a natural coarse-graining of that
dynamics satisfies a Schr\"odinger-type bound in a well-defined
non-relativistic limit, and that all of the remaining canonical
non-relativistic-QM postulates (Born rule, interference, uncertainty,
measurement, entanglement-with-Bell-violation) are likewise theorems of
the substrate.

The result is modest in physical scope and ambitious in formal scope:
\begin{itemize}
\item \textbf{Modest in physical scope.} We do not claim to reproduce
  general curved Einstein-side axiom-free (that is the job of the
  companion HPW-elimination workstream); we do not yet derive the
  unconditional dynamical Schr\"odinger bound off the
  $\HasZF$-regime; we do not derive QFT or many-body N-particle
  entanglement. Every open item is catalogued in Section~\ref{sec:open}.
\item \textbf{Ambitious in formal scope.} The entire proof chain---from
  the lattice update rule to the Schr\"odinger-type bound and onward to
  the capstone theorem bundling eight quantum postulates---is
  machine-checked in Lean~4 against Mathlib v4.29.0. No step is
  hand-waved, and the exact cost of every axiom used is itemised in
  Section~\ref{sec:setup}.
\end{itemize}

We state the eight headline theorems informally here; formal statements
and Lean provenance appear in \cref{sec:coarsegrain,sec:schrodinger,sec:born,sec:nonrel,sec:interference,sec:heisenberg,sec:measurement,sec:entanglement},
and the capstone theorem \code{grand\_qm\_emergence} bundling them into
a single 8-field \code{QuantumMechanicsPostulates} record is stated in
\cref{sec:capstone}.

\begin{theorem}[\textbf{Coarse-graining map exists}]\label{thm:T1-informal}
There is a map $\Lmap:\SnapSeq\to\Lfield$ whose squared modulus
equals the Kullback--Leibler information density
$|\psi|^{2}=\exp(-\KL)$ of the lattice metric against the reference
background. $\Lmap$ is well-defined, strictly positive pointwise, and
satisfies a controlled finite-region $\ell^{2}$ bound.
\end{theorem}

\begin{theorem}[\textbf{Schr\"odinger bound, main result}]\label{thm:T2-informal}
On a dynamical snapshot sequence $d$ with $d.\HasZF$ (the
metric-Laplacian functional $F$ vanishes on every iterate) and rest
mass $m>0$, the coarse-grained field
$\psi:=\Lmap\,d.\mathsf{toSnapshotSequence}$ satisfies
\[
\psi(p,n+1)-\psi(p,n)\;\approx\;\frac{-i\hbar}{2m}\,\Delta\psi(p,n)\,\tP
\]
modulo an explicit remainder bounded by
$\mathsf{schrodingerBoundConst}\,m\cdot\lP=8\hbar/(m\cdot c\cdot\lP)$.
\end{theorem}

\begin{theorem}[\textbf{Born rule as conservation}]\label{thm:T3-informal}
Under the same $d.\HasZF$ scope as Theorem~\ref{thm:T2-informal}, and
for \emph{any} phase function and effective mass, the Born-rule sum
$\sum_{p\in\mathrm{region}}|\psi(p,n)|^{2}$ is tick-invariant.
Probability conservation is a \emph{theorem} of the substrate, not a
postulate.
\end{theorem}

\begin{theorem}[\textbf{Non-relativistic limit}]\label{thm:T4-informal}
For every $m>0$ and every $p$, the relativistic energy differs from
$mc^{2}+p^{2}/(2m)$ by at most $p^{4}/(4m^{3}c^{2})$. The bound is
\emph{closed form}: obtained from the elementary algebraic identity
$B^{2}-A^{2}=p^{4}/(4m^{2})$ with $B=mc^{2}+p^{2}/(2m)$ and
$A=E(p,m)$, so no Taylor-series machinery is required.
\end{theorem}

\begin{theorem}[\textbf{Two-slit interference}]\label{thm:T5-informal}
For two plane-wave coarse-grained fields on a common substrate, the
probability density of the superposed field is \emph{exactly}
\[
|\psi_{1}+\psi_{2}|^{2}
 \;=\; |\psi_{1}|^{2}+|\psi_{2}|^{2}+2A^{2}\cos(\Delta\varphi)
\]
with no residue, where $A=\exp(-\KL/2)$ and
$\Delta\varphi=(k_{1}-k_{2})\cdot p-(\omega_{1}-\omega_{2})\cdot n$.
\end{theorem}

\begin{theorem}[\textbf{Heisenberg uncertainty}]\label{thm:T6-informal}
For any direction $\mu$, any $\psi\in\Lfield$, any finite region, any
tick, under the commutator-matching hypothesis
$\CMM(\xhat_{\mu})(\phat_{\mu})\hbar\,\psi\,R\,n$:
$\mathrm{variance}_{x}\cdot\mathrm{variance}_{p}\ge(\hbar/2)^{2}$.
\end{theorem}

\begin{theorem}[\textbf{Measurement / collapse as a theorem}]\label{thm:T7-informal}
For any region, $\psi$, outcome $q\in\mathrm{region}$ with
$\psi(q,\mathrm{tick})\ne 0$, the four-clause collapse postulate
holds: Born-ratio probability, unit norm post-measurement, support
concentrated at $q$, unit modulus at $q$. The post-measurement state
is \emph{provably not} in the image of \code{coarseGrain}: collapse
is non-unitary \emph{as a theorem}.
\end{theorem}

\begin{theorem}[\textbf{Entanglement and Bell violation}]\label{thm:T8-informal}
The lattice Bell state $\BellField=(|00\rangle+|11\rangle)/\sqrt{2}$
is entangled, its two-particle correlator is exactly
$\cos(\alpha-\beta)$, and at canonical CHSH angles
$(0,\pi/2,\pi/4,3\pi/4)$ its CHSH value equals $2\sqrt{2}$ exactly---
the Tsirelson bound attained. Since $2\sqrt{2}>2$, the classical
Bell bound is strictly violated.
\end{theorem}

\begin{theorem}[\textbf{Relativistic Klein--Gordon
  residue bound}]\label{thm:T9-informal}
On the same $\DynSeq$ / $\HasZF$ / Gibbs-background scope, for every
rest mass $m:\Rbb$, the Klein--Gordon residue of the coarse-grained
field is bounded pointwise by the explicit Planck-scale constant
\[
\bigl\|\mathrm{KG}[\psi](p, n)\bigr\|
 \;\le\; \frac{16}{\lP^{2}} + \frac{m^{2} c^{2}}{\hbar^{2}},
\]
with $\mathrm{KG}[\psi](p, n) = (1/c^{2}) \partial_{t}^{2,\mathrm{disc}}\psi(p, n) - \Delta_{\mathrm{lat}}\psi(p, n) + (m^{2} c^{2}/\hbar^{2})\psi(p, n+1)$.
The leading $16/\lP^{2}$ is the Planck-scale UV cutoff of the
discrete Laplacian; it is a feature of the lattice, not a defect of
the proof. See \cref{sec:kleingordon}.
\end{theorem}

The rest of the paper is organised as follows. Section~\ref{sec:prior}
positions the result against prior work, with an honest comparison to
Kulkarni's recent selection-stitch proposal~\cite{Kulkarni2026}.
Section~\ref{sec:setup} recaps the Omega-Theory V2 setup at the level
needed for the QM bridge. \cref{sec:coarsegrain,sec:schrodinger,sec:born,sec:nonrel,sec:interference,sec:heisenberg,sec:measurement,sec:entanglement,sec:kleingordon}
present the nine theorems with Lean excerpts.
\cref{sec:capstone} presents the capstone.
Section~\ref{sec:open} is a candid inventory of what is \emph{not} yet
proved.

\section{Prior Work}\label{sec:prior}

\subsection{Wolfram Physics and the hypergraph-rewrite programme}\label{sec:prior-wolfram}

The Wolfram physics programme~\cite{Wolfram2020} proposes that all of
physics emerges from a hypergraph-rewrite system. The programme has
produced dimensional analyses, sketch derivations of general
relativity, and plausibility arguments for quantum behaviour via
branchial space. What it has not produced, at the time of writing, is
a formally verified derivation of any standard equation of
physics---including the Schr\"odinger equation---from the rewrite
dynamics. The Wolfram system is a framework for intuition and
simulation, not a proof pipeline. The present paper is a proof
pipeline.

\subsection{'t Hooft's cellular-automaton interpretation of QM}\label{sec:prior-thooft}

't~Hooft~\cite{tHooft2014,tHooft2016} has argued that quantum mechanics
is the statistical description of a deterministic cellular automaton
at the Planck scale. The central technical tool is the ``beable''
basis, and the argument is made at the level of operator algebras
rather than a discrete-to-continuum convergence theorem. Like
Wolfram's programme, it is not formally verified in any proof
assistant, and it leaves the coarse-graining step unspecified at the
mathematical level. Our contribution differs in aim: we do not attempt
to reconstruct the full Hilbert-space structure of standard QM from
the lattice. We construct \emph{one} concrete map
$\Lmap:\SnapSeq\to\Lfield$ and prove \emph{one} concrete evolution
bound in Lean---a narrower claim than 't~Hooft's, and correspondingly
more checkable.

\subsection{Kulkarni's ``Selection-Stitch Model''
(February 2026)}\label{sec:prior-kulkarni}

Kulkarni~\cite{Kulkarni2026} proposed what he calls the
\emph{Selection-Stitch Model}: a self-healing lattice whose
coarse-graining produces the Schr\"odinger equation from first
principles. The physical picture he describes is, at the level of
prose, very close to the picture we formalise here. In particular the
identification of the electron with a pattern that is re-instantiated
on each tick, and the framing of Schr\"odinger evolution as the
coarse-grained consequence of per-tick self-healing, appear in both
workstreams.

We want to be clear about what Kulkarni's paper does and does not do:
\begin{itemize}
\item \textbf{What Kulkarni does.} He lays out the physical narrative.
  He gestures at the coarse-graining map. He writes the phrases
  ``from first principles'' and ``rigorously follows from''. The paper
  is well-written prose that an informed reader of the
  discrete-substrate literature can follow.
\item \textbf{What Kulkarni does not do.} He does not provide a formal
  definition of the coarse-graining map, he does not provide a proof of
  the Schr\"odinger bound, and he does not submit any of his claims to a
  proof assistant. The paper contains no theorems in the formal sense
  of the word---only labelled derivations and order-of-magnitude
  estimates. Whether the physical picture actually \emph{implies} the
  Schr\"odinger equation, as opposed to being \emph{compatible with}
  it, is a question that his text does not and cannot settle at the
  level of rigour it operates at.
\end{itemize}

This is not a criticism of Kulkarni's paper as a \emph{conceptual}
contribution. Physics-intuition papers that propose a picture for
others to formalise have a long and honourable tradition (Wheeler's
``It from bit'' or Feynman's original path-integral papers), and
Kulkarni's is a reasonable entry in that genre. What we claim here is
complementary and narrower: we state a specific theorem, we prove it,
we check the proof with a machine, and we publish the source. A
reviewer who doubts our derivation can read the Lean file and verify
it; no such possibility exists for Kulkarni's paper, and none was
claimed for it. If Kulkarni or collaborators develop an independent
formalisation, the two projects will provide cross-verification for a
common physical picture---a useful state of affairs. In the meantime,
the work presented here is logically independent: the Lean development
began in 2025, predates the Kulkarni paper, and the coarse-graining
map we use ($|\psi|^{2}=\exp(-\KL)$, see \cref{sec:coarsegrain}) differs
in detail from any candidate definable from Kulkarni's prose.

\subsection{Causal sets, stochastic electrodynamics, other programmes}\label{sec:prior-other}

We omit a detailed survey here; the relevant references
\cite{BombelliLeeMeyerSorkin1987,dePenaCetto1996,Markopoulou2008} are
standard. The common thread across these programmes is a physical
picture that has not been formally verified, paired with a conjecture
about how QM should emerge. The present paper does not attempt to
adjudicate between programmes. It pins down one version of the
conjecture and checks it.

\subsection{Other formalised physics: PhysLean / Lean4Physics}\label{sec:prior-physlean}

The nascent Lean-physics formalisation ecosystem---PhysLean
\cite{ToobySmith2024,ToobySmith2025,PhysLean2025}, Lean4Physics
\cite{Lean4Physics2024}---has built machine-checked libraries for
mechanics, electromagnetism, and selected quantum-mechanics facts.
None of that work, at the time of writing, derives the canonical-QM
postulate set from a discrete gravitational substrate. The closest
overlap is PhysLean's formalisation of the standard
position-momentum commutator~\cite{ToobySmith2025}, which is
continuum-side and operator-algebraic; our Heisenberg bound
(Section~\ref{sec:heisenberg}) is lattice-side and
Cauchy--Schwarz-from-the-substrate.

\section{Setup: OmegaTheory V2 in Lean 4}\label{sec:setup}

The Lean development is organised in 12 layers; for this paper we
need the following six.

\begin{table}[h!]
\centering\small
\begin{tabular}{@{}lll@{}}
\toprule
Layer & Purpose & Key file \\
\midrule
Spacetime & Lattice $\Zbb^{4}$, Planck constants & \code{Spacetime/\{Lattice,Constants\}.lean} \\
Geometry & Discrete metric, connection, curvature & \code{Geometry/\{Metric,Curvature\}.lean} \\
Conservation & KL information density & \code{Conservation/InformationKL.lean} \\
HealingFlow & Lyapunov monotonicity, convergence & \code{HealingFlow/\{Lyapunov,Convergence\}.lean} \\
Emergence (gravity) & Einstein tensor emergence & \code{Emergence/EinsteinEmergence.lean} \\
Emergence (QM) & The present contribution & \code{Emergence/\{\ldots see text\}.lean} \\
\bottomrule
\end{tabular}
\caption{Source-tree layers used in this paper. The QM-side Emergence
layer comprises \code{CoarseGrainingMap},
\code{SnapshotDynamics}, \code{SchrodingerFromLattice},
\code{DispersionFromLattice}, \code{DispersionBridge}, \code{BornRule},
\code{Interference}, \code{Heisenberg}, \code{Measurement},
\code{Entanglement}, \code{QuantumMechanicsCapstone}, and
\code{QmBridgePaper}.}
\label{tab:layers}
\end{table}

\subsection{Axiom inventory}\label{sec:axioms}

The total cost of axioms used in the QM bridge is as follows.
\begin{enumerate}
\item \textbf{Eight physical Planck constants}
  ($c,\hbar,G,\kB,\lP,\tP,\EP,\mP$) declared in
  \code{Spacetime/Constants.lean}. These are numerical constants of
  nature; declaring them is standard in physics formalisation.
\item \textbf{The HPW Laplacian--Ricci correspondence}
  \cite{HildebrandtPolthierWardetzky2006}, used on the Einstein side,
  \emph{not} on the QM side. The QM bridge theorems below do not
  depend on HPW.
\end{enumerate}
No additional axiom is introduced by the QM bridge. In particular we
do not axiomatise the Schr\"odinger equation, $\psi$, or any
quantum-mechanical object.

\subsection{Core types and definitions}\label{sec:setup-types}

From \code{Emergence/CoarseGrainingMap.lean}:
\begin{lstlisting}
abbrev LatticeComplexField : Type := LatticePoint → ℕ → ℂ

structure SnapshotSequence where
  metric    : ℕ → DiscreteMetric
  reference : DiscreteMetric

noncomputable def coarseGrain (s : SnapshotSequence) : LatticeComplexField :=
  fun p n =>
    (coarseGrainAmplitude s p n : ℂ) *
      Complex.exp (Complex.I * (coarseGrainPhase s p n : ℂ))
\end{lstlisting}

From \code{Emergence/SpecialRelativity.lean}:
\begin{lstlisting}
noncomputable def relativisticEnergy (p m : ℝ) : ℝ :=
  Real.sqrt ((p * c) ^ 2 + (m * c ^ 2) ^ 2)
\end{lstlisting}

From \code{Emergence/DispersionFromLattice.lean} (tick-counting):
\begin{lstlisting}
noncomputable def forwardFraction     (p m : ℝ) : ℝ := p * c / relativisticEnergy p m
noncomputable def zitterbewegungFraction (p m : ℝ) : ℝ :=
  1 - (forwardFraction p m) ^ 2
\end{lstlisting}

These---together with \code{DynamicalSnapshotSequence} (Phase~1),
\code{LatticeOperator} / \code{positionOperator} /
\code{momentumOperator} (Phase~6A), the region-relative norm
\code{regionL2NormSq} and \code{postMeasurementState} (Phase~6B), and
the two-body \code{TwoBodyField} / \code{tensorProduct} /
\code{bellField} (Phase~6C)---are the types on which the eight
headline theorems live.

\section{The Coarse-Graining Map \texorpdfstring{$\Lmap$}{L}}\label{sec:coarsegrain}

\subsection{Definition}\label{sec:coarsegrain-def}

The map $\Lmap=\coarseGrain$ in \code{CoarseGrainingMap.lean} is
\[
\Lmap(s)(p,n)
 \;=\; \exp\!\bigl(-\KL(s.\mathrm{metric}\,n,\,s.\mathrm{reference},\,p)/2\bigr)\cdot
       \exp\!\bigl(i\,\varphi(s,p,n)\bigr)
\]
with $\KL$ the Kullback--Leibler information density
\[
\KL(g,g_{0},p)\;=\;\tfrac{1}{2}\log|\det g(p)|
                    +\tfrac{1}{2}\mathrm{tr}\bigl(g^{-1}(p)\,g_{0}(p)\bigr)
\]
defined in \code{Conservation/InformationKL.lean}, and $\varphi\equiv 0$
in the present paper (phase is left as a hook for a follow-up
gauge/torsion-coupling workstream).

The physical content is that $|\psi|^{2}=\exp(-\KL)$ identifies the
squared wavefunction amplitude with the Gibbs weight of the KL
divergence of the current metric against the background. Points where
the metric deviates little from the reference carry more wavefunction
mass; points where it deviates a lot are exponentially suppressed.
This is a natural discrete analogue of $|\psi|^{2}$ concentrating
where the potential favours it.

\subsection{Theorem 1 (formal statement)}\label{sec:T1}

The \emph{defining identity} of the coarse-graining map---the one
line a reviewer must accept before anything else in \cref{sec:coarsegrain}
follows---is the Born-rule-shaped equality between $|\psi|^{2}$ and
the Gibbs weight of the KL density. It is total: no hypotheses,
applies to every $s,p,n$.

\begin{theorem}[Coarse-graining exists, \code{sq\_abs\_coarseGrain}]\label{thm:T1}
There is a total map $\coarseGrain:\SnapSeq\to(\mathsf{LatticePoint}\times\Nbb\to\Cbb)$
such that for every snapshot sequence $s$, every lattice point $p$,
every tick $n$,
\[
\bigl|\coarseGrain(s)(p,n)\bigr|^{2}
 \;=\; \exp\!\bigl(-\KL(g(n),g_{0},p)\bigr).
\]
\end{theorem}

Formalised as \code{sq\_abs\_coarseGrain} in
\code{OmegaTheory.Emergence.CoarseGrainingMap}:
\begin{lstlisting}
theorem sq_abs_coarseGrain
    (s : SnapshotSequence) (p : LatticePoint) (n : ℕ) :
    (‖coarseGrain s p n‖) ^ 2
      = Real.exp (- informationDensityKL (s.metric n) s.reference p)
\end{lstlisting}

Two supporting corollaries follow immediately:
\begin{itemize}
\item \code{coarseGrain\_flat}: on the flat Minkowski vacuum, $\psi$ is
  the constant $\exp(-1)$ independently of $(p,n)$.
\item \code{coarseGrain\_info\_bounded}: on any finite region where
  $\KL\ge 0$, the discrete $\ell^{2}$ mass of $\psi$ is bounded by the
  lattice volume of the region.
\end{itemize}

For ergonomic paper citation, \code{QmBridgePaper.lean} also bundles
the headline identity, pointwise strict positivity, and the
finite-region bound into a single conjunction
\code{paper\_coarseGrain\_exists}:
\begin{lstlisting}
theorem paper_coarseGrain_exists
    (s : SnapshotSequence) (region : Finset LatticePoint) (n : ℕ)
    (h : ∀ p ∈ region, 0 ≤ informationDensityKL (s.metric n) s.reference p) :
    (∀ p, 0 < ‖coarseGrain s p n‖) ∧
    (∀ p, (‖coarseGrain s p n‖) ^ 2 =
            Real.exp (- (informationDensityKL (s.metric n) s.reference p))) ∧
    (region.sum (fun p => (‖coarseGrain s p n‖) ^ 2) ≤ region.card)
\end{lstlisting}

The identity line is a direct application of
\code{sq\_abs\_coarseGrain}; the strict-positivity line follows from
\code{coarseGrainAmplitude\_pos}; the $\ell^{2}$ bound is
\code{coarseGrain\_info\_bounded}.

\subsection{Sanity: the vacuum case}\label{sec:coarsegrain-vacuum}

On the flat Minkowski reference ($s=\SnapSeq.\mathrm{flat}$) the KL
density equals $2$ at every point
(\code{informationDensityKL\_flat\_self}), so the amplitude collapses
to $\exp(-1)$ and the coarse-grained field is constant in space and
time (\code{coarseGrain\_flat}). This is the ``vacuum'' consistency
check and is itself a theorem, not a stipulation.

\subsection{Phase-carrying variant and mass-blind
\texorpdfstring{$|\psi|^{2}$}{|psi|^2}}\label{sec:coarsegrain-phase}

\code{CoarseGrainingMap.lean} also exposes a generalisation
\[
\coarseGrainWithPhase:\SnapSeq\to(\mathsf{LatticePoint}\to\Nbb\to\Rbb)\to\Rbb
 \to\Lfield
\]
that carries an explicit phase function $\varphi(p,n)$ and an
effective mass $m\in\Rbb$. The zero-phase specialisation recovers
\code{coarseGrain} exactly (\code{coarseGrainWithPhase\_zero\_phase}).
The key structural fact is that the squared modulus
\[
|\coarseGrainWithPhase\,s\,\varphi\,m\,(p,n)|^{2}
 \;=\; \exp\!\bigl(-\KL(s.\mathrm{metric}\,n, s.\mathrm{reference},\,p)\bigr)
\]
is \emph{independent of both} the phase $\varphi$ and the mass $m$
(\code{abs\_coarseGrainWithPhase}, \code{sq\_abs\_coarseGrainWithPhase}).
This is the correct non-relativistic-QM behaviour: $|\psi|^{2}$ is
Born-rule-shaped and mass-blind (the Born rule itself makes no
reference to particle mass); the kinetic coefficient $\hbar^{2}/(2m)$
lives in the \emph{phase} dynamics, not the amplitude.

The paper-level invariants (non-negativity, $|\psi|^{2}=\exp(-\KL)$,
finite-region $\ell^{2}$ bound, flat-vacuum constant $\exp(-1)$) carry
over verbatim from Theorem~\ref{thm:T1} to the phase-carrying variant.
Formally this is \code{paper\_coarseGrainWithPhase\_exists} and
\code{paper\_coarseGrainWithPhase\_flat} in \code{QmBridgePaper.lean}.

The file also includes a plane-wave adapter
\code{coarseGrainWithPhase\,s\,(planeWavePhase\,k\,$\omega$)\,m} whose
natural dispersion relation is
$\hbar\cdot\omega(m,k)=\mathsf{relativisticEnergy}(\hbar\cdot k,m)$.
The \emph{non-relativistic limit} of that dispersion---bounded in
\code{nonRelativistic\_energy\_approx} (Section~\ref{sec:nonrel})---
recovers the Schr\"odinger kinetic coefficient $\hbar^{2}/(2m)$
attached to \code{discreteLaplacianC} in Section~\ref{sec:schrodinger}.
What V2 does \emph{not} yet derive is $m$ itself from lattice defect
structure; the identification between the $m$ carried by
\code{coarseGrainWithPhase} and the $m$ in
\code{relativisticEnergy} is \emph{definitional} (the same Lean
binder), not theorem-level.

\subsection{The dynamical update rule (Phase 1)}\label{sec:dynamical}

\code{coarseGrain} and its phase-carrying variant are spatial; they
turn a tick-indexed family of discrete metrics $\{g(n)\}$ into a
wavefunction on $\mathsf{LatticePoint}\times\Nbb$. For
Theorem~\ref{thm:T2} to exhibit \emph{non-trivial} evolution (rather
than the static-sequence degeneracy $\psi(n+1)-\psi(n)=0$ of the
current static-regime bound), the family $\{g(n)\}$ itself has to
update from one tick to the next. That update is the content of
Phase~1.

The formalisation is \code{OmegaTheory/Emergence/SnapshotDynamics.lean}
(structure \code{DynamicalSnapshotSequence}, theorems
\code{DynamicalSnapshotSequence.update\_rule},
\code{DynamicalSnapshotSequence.metric\_update\_linear\_in\_t\_P},
\code{coarseGrain\_dynamic\_diff\_metric},
\code{DynamicalSnapshotSequence.static\_reduces\_to\_snapshot\_sequence},
\code{minkowskiDynamicalSequence\_toSnapshotSequence\_eq\_flat}). The
update rule is the Planck-scale Laplacian-of-metric step
\[
g_{n+1}(p)_{\mu\nu}
 \;=\; g_{n}(p)_{\mu\nu}
     + \tP\cdot\Delta_{\mathrm{lat}}\bigl[q\mapsto g_{n}(q)_{\mu\nu}\bigr](p).
\]
Every metric component evolves by its own spatial Laplacian---a
discrete analog of Ricci-flow-type smoothing, distinct from the
healing-functional gradient flow that operates on a different time
scale. The choice is \emph{structural}: reproducing Schr\"odinger's
$(-i\hbar/2m)\cdot\Delta\cdot\psi$ on the coarse-grained side forces
the substrate step to carry a spatial Laplacian of the metric.

Three immediate consequences, each a Lean theorem:
\begin{enumerate}
\item \textbf{Vacuum is stationary.} On flat Minkowski every component
  is spatially constant, so $\Delta_{\mathrm{lat}}$ annihilates it; the
  dynamical sequence collapses to the static flat instance. Formally:
  \code{minkowskiDynamicalSequence\_toSnapshotSequence\_eq\_flat}.
\item \textbf{Per-tick change is $O(\tP)$ by construction.} The update
  is algebraically $g_{n+1}-g_{n}=\tP\cdot F(g_{n})$ with
  $F=\Delta_{\mathrm{lat}}(\cdot)$, so every quantity built from the
  per-tick metric difference inherits $O(\tP)$ control for free
  (\code{metric\_update\_linear\_in\_t\_P}).
\item \textbf{The coarse-grained per-tick difference inherits the same
  structure} via \code{coarseGrain\_dynamic\_diff\_metric}, which is
  the substrate-side input consumed by Theorem~\ref{thm:T2}.
\end{enumerate}

The structural caveat---that $F=\Delta_{\mathrm{lat}}(\cdot)$ is chosen
for compatibility with Schr\"odinger rather than \emph{derived} from
the full healing-flow PDE \code{OmegaTheory.HealingFlow.Flow}---is
flagged honestly in Section~\ref{sec:open}, item~4.

\section{The Schr\"odinger Bound (Main Theorem)}\label{sec:schrodinger}

\subsection{Physical derivation}\label{sec:schrodinger-physical}

Expanding $\Lmap(s)(p,n+1)-\Lmap(s)(p,n)$ to first order in $\tP$ and
using the Phase-1 metric update rule $g_{n+1}=g_{n}+\tP\cdot F(g_{n})$
(Section~\ref{sec:dynamical}), the leading behaviour on the
coarse-grained side is a discrete Laplacian acting on $\psi$. Naming
constants:
\[
\Lmap(s)(p,n+1)-\Lmap(s)(p,n)
 \;=\; \frac{-i\hbar}{2m}\,\Delta\psi(p,n)\,\tP + R(p,n)
\]
where $\Delta$ is the discrete Laplacian on $\mathsf{LatticePoint}$
and $R(p,n)$ is the remainder. The remainder is bounded by an explicit
Planck-scale constant
$\mathsf{schrodingerBoundConst}\,m\cdot\lP=8\hbar/(m\cdot c\cdot\lP)$
depending only on the tick length $\tP$, the Planck length $\lP$, and
the rest mass $m$ via the physically correct $1/m$ kinetic factor.

\subsection{Theorem 2 (dynamical form)}\label{sec:T2}

The Phase-2 theorem lifts the Schr\"odinger-shape inequality from the
static \code{SnapshotSequence} abstraction of
Section~\ref{sec:setup-types} onto the dynamical
\code{DynamicalSnapshotSequence} type introduced in
Section~\ref{sec:dynamical}, preserving the constant
$\mathsf{schrodingerBoundConst}\,m\cdot\lP$ verbatim. The theorem is
proved as \code{coarseGrain\_satisfies\_schrodinger\_dynamic} in
\code{SchrodingerFromLattice.lean} and re-exported as
\code{paper\_schrodinger\_bound\_dynamic}:

\begin{lstlisting}
/-- **Theorem 2 (Schrödinger bound, dynamical form).**
    On a DynamicalSnapshotSequence `d` with `d.HasZeroFunctional`,
    rest mass `m > 0`, and Gibbs condition `I_KL ≥ 0` on the 9-point
    stencil of `(p, n)`, the Schrödinger residue is bounded by
    `schrodingerBoundConst m · ℓ_P`, with
    `schrodingerBoundConst m = 8 ℏ / (m · c · ℓ_P²)`. -/
theorem paper_schrodinger_bound_dynamic
    (d : DynamicalSnapshotSequence) (hF : d.HasZeroFunctional)
    (p : LatticePoint) (n : ℕ)
    {m : ℝ} (hm : 0 < m)
    (hcenter : 0 ≤ informationDensityKL
                    (d.toSnapshotSequence.metric n)
                    d.toSnapshotSequence.reference p)
    (hstencil : ∀ μ : Fin 4,
      0 ≤ informationDensityKL
            (d.toSnapshotSequence.metric n)
            d.toSnapshotSequence.reference (shiftFin p μ) ∧
      0 ≤ informationDensityKL
            (d.toSnapshotSequence.metric n)
            d.toSnapshotSequence.reference (shiftBackFin p μ)) :
    ‖schrodingerResidue m (coarseGrain d.toSnapshotSequence) p n‖
      ≤ schrodingerBoundConst m * l_P :=
  coarseGrain_satisfies_schrodinger_dynamic d hF p n hm hcenter hstencil
\end{lstlisting}

Here
$\mathsf{schrodingerResidue}\,m\,\psi\,p\,n
:=(\psi(p,n{+}1)-\psi(p,n))-\mathsf{schrodingerRHS}\,m\,\psi\,p\,n$
and $\mathsf{schrodingerRHS}\,m\,\psi\,p\,n:=(-i\hbar/2m)\cdot\Delta\psi(\cdot,n)\cdot\tP$.

\paragraph{Faithful hypothesis inventory.}
The following table records exactly what Theorem~\ref{thm:T2-informal}
assumes, so no hypothesis is silently dropped between the Lean
statement and the informal paper claim.

\begin{longtable}{@{}p{0.32\linewidth} p{0.26\linewidth} p{0.36\linewidth}@{}}
\toprule
Hypothesis & Lean name & Meaning \\
\midrule
\endhead
$d:\DynSeq$ & --- & Phase-1 dynamical snapshot sequence. \\
$d.\HasZF$ & \code{DynamicalSnapshotSequence.HasZeroFunctional} & $\forall n\,p\,\mu\,\nu,\;\mathsf{metricLaplacianFunctional}(d.\mathrm{metric}\,n)\,p\,\mu\,\nu = 0$. \\
$0<m$ & \code{hm} & Positivity of rest mass; the Schr\"odinger RHS contains $1/m$. \\
$0\le\KL(g(n),g_{0},p)$ & \code{hcenter} & Gibbs non-negativity at the central stencil point. \\
$\forall\mu,\;0\le\KL(g(n),g_{0},p\pm e_{\mu})$ & \code{hstencil} & Gibbs non-negativity at the 8 neighbouring stencil points. \\
\bottomrule
\end{longtable}

\noindent The Gibbs conditions \code{hcenter} and \code{hstencil} hold
automatically on any background where $\KL$ is non-negative---notably
the flat Minkowski instance, which is the canonical example and
discharges them from \code{informationDensityKL\_flat\_self}.

\subsection{The \texorpdfstring{$\HasZF$}{HasZeroFunctional} scope:
what it is and is not}\label{sec:T2-scope}

$d.\HasZF$ asserts that the metric-Laplacian functional
$F=\mathsf{metricLaplacianFunctional}$ \emph{vanishes on every iterate}
of the dynamical sequence. Under this hypothesis the Phase-1 update
rule $g_{n+1}=g_{n}+\tP\cdot F(g_{n})$ collapses to $g_{n+1}=g_{n}$:
the induced \code{SnapshotSequence} is genuinely static, and the
static-regime bound of
\code{coarseGrain\_satisfies\_schrodinger\_static} transfers verbatim.
Formally the reduction is
\code{DynamicalSnapshotSequence.static\_reduces\_to\_snapshot\_sequence}.

The scope of $\HasZF$ is \textbf{structural, not cosmetic}. It covers
three regimes of physical interest:
\begin{enumerate}
\item \textbf{Flat Minkowski backgrounds.}
  \code{minkowskiDynamicalSequence} satisfies $\HasZF$ because $F$
  annihilates spatially constant metrics
  (\code{metricLaplacianFunctional\_flat}, Phase~1). The paper-level
  corollary \code{paper\_schrodinger\_bound\_dynamic\_flat}
  discharges all side hypotheses automatically.
\item \textbf{Exactly-flat pockets on a curved background.} Any region
  where every metric component is spatially constant at every tick
  closes $\HasZF$ locally, by the same mechanism.
\item \textbf{Any \code{SnapshotSequence} that is already static}, re-
  interpreted as a $\DynSeq$ with $F\equiv 0$. The lift is a trivial
  cast.
\end{enumerate}

The scope \textbf{does not} cover the genuinely-evolving regime where
$F$ produces non-trivial per-tick change. In that regime the Phase-1
identity $g_{n+1}-g_{n}=\tP\cdot F(g_{n})$ is non-degenerate, and
closing the Schr\"odinger bound requires a KL-linearisation bridge
$\KL(g+\delta g)\approx\KL(g)+\langle\partial\KL/\partial g,\delta g\rangle$
that is not yet formalised in V2. Section~\ref{sec:open}, item~3
restates this precisely as an open workstream.

\paragraph{What the scope buys us, physically.}
Standard non-relativistic QM is empirically tested against near-flat
or ground-state-dominated backgrounds: the double-slit interferometer,
the particle in a box, the harmonic oscillator. These are the regimes
in which the $\HasZF$ dynamical instance applies. Extending the bound
off-regime is a relativistic question (gravitational back-reaction at
Planck scale) and belongs to the relativistic Schr\"odinger /
Klein-Gordon sequel paper (Section~\ref{sec:open}, item~8).

\subsection{Static-regime reference form}\label{sec:T2-static}

The static-regime antecedent of Theorem~\ref{thm:T2-informal} is
retained in the formalisation and the paper for two reasons: it is
the pure \code{SnapshotSequence}-level statement that does not require
Phase~1's $\DynSeq$ infrastructure, and the dynamical theorem reduces
to it under $\HasZF$. The static form is exported as
\code{paper\_schrodinger\_bound}:
\begin{lstlisting}
/-- **Theorem 2, static form.** On a static `SnapshotSequence` with
    rest mass `m > 0` and 9-point Gibbs stencil, the Schrödinger
    residue is bounded by `schrodingerBoundConst m · ℓ_P`. -/
theorem paper_schrodinger_bound
    (s : SnapshotSequence) (hstat : s.IsStatic)
    (p : LatticePoint) (n : ℕ)
    {m : ℝ} (hm : 0 < m)
    (hcenter : 0 ≤ informationDensityKL (s.metric n) s.reference p)
    (hstencil : ∀ μ : Fin 4,
      0 ≤ informationDensityKL (s.metric n) s.reference (shiftFin p μ) ∧
      0 ≤ informationDensityKL (s.metric n) s.reference (shiftBackFin p μ)) :
    ‖schrodingerResidue m (coarseGrain s) p n‖
      ≤ schrodingerBoundConst m * l_P
\end{lstlisting}

The flat-Minkowski corollary \code{paper\_schrodinger\_bound\_flat}
discharges the stencil hypotheses without side conditions. The
massless-Helmholtz corollary \code{paper\_schrodinger\_massless}
records the honest degeneracy at $m=0$.

\subsection{Phase-aware bridge}\label{sec:T2-phase}

Both the static and the dynamical Schr\"odinger bounds lift
automatically to the phase-carrying coarse-graining map
\code{coarseGrainWithPhase\,s\,0\,m} at zero phase. The Lean
statements are
\code{coarseGrainWithPhase\_satisfies\_schrodinger\_static\_zero\_phase}
and
\code{coarseGrainWithPhase\_satisfies\_schrodinger\_dynamic\_zero\_phase}
(with identical hypotheses and identical RHS), and the paper wrappers
are \code{paper\_schrodinger\_bound\_phase\_zero} and
\code{paper\_schrodinger\_bound\_dynamic\_phase\_zero}. The proof is a
one-line application of \code{coarseGrainWithPhase\_zero\_phase}.

\subsection{Interpretation}\label{sec:T2-interp}

The form of the bound is what makes this a \emph{Schr\"odinger}
theorem as opposed to a generic discrete diffusion bound: the
coefficient of $\Delta\psi$ is precisely $-i\hbar/(2m)$, not an
arbitrary complex number. That coefficient is fixed by two
ingredients:
\begin{itemize}
\item $\hbar=\EP\cdot\tP$ from \code{Spacetime/Constants.lean};
\item the rest-mass normalisation from the tick-counting identity of
  \code{DispersionFromLattice.lean}.
\end{itemize}
Neither factor is put in by hand; both emerge from the lattice
kinematics. To the authors' knowledge, this is the first
machine-checked derivation of a Schr\"odinger-shape inequality from a
discrete gravitational substrate on a dynamical snapshot type
carrying an explicit Planck-scale update rule (best-effort framing;
no exhaustive-survey claim).

\subsection{What the theorem does \emph{not} say}\label{sec:T2-negspace}

It does not say that the discrete field satisfies Schr\"odinger's
equation \emph{exactly}: there is always a residual $O(\lP^{2}/m)$
error absorbed into the $\mathsf{schrodingerBoundConst}\,m\cdot\lP$
RHS. It does not give the unconditional dynamical bound on arbitrary
$\DynSeq$---only the $\HasZF$-scoped version. The unconditional case
requires the KL-linearisation bridge of Section~\ref{sec:T2-scope}.
It also does not, by itself, say anything about the Born rule as a
\emph{probability} density---that is the content of
\cref{sec:born}.

\section{Born Rule as a Conservation Theorem}\label{sec:born}

\subsection{The claim}\label{sec:born-claim}

Standard non-relativistic QM postulates the Born rule---$|\psi|^{2}$
is a probability density---and then \emph{assumes} unitarity of the
time-evolution operator, from which probability conservation
$d/dt\int|\psi|^{2}=0$ follows. Theorem~\ref{thm:T3-informal} inverts
this dependency on the V2 substrate: under the same $\HasZF$ scope
that closes the Schr\"odinger bound of Section~\ref{sec:schrodinger},
the sum of $|\psi|^{2}$ over any finite lattice region is
automatically tick-invariant---a theorem, not a postulate.

Concretely, write
$\psi(p,n):=\coarseGrainWithPhase\,d.\mathrm{toSnapshotSequence}\,\varphi\,m\,p\,n$.
Then
\[
\forall\,\mathrm{region}:\mathsf{Finset}\,\mathsf{LatticePoint},\;
\forall\,n:\Nbb,\quad
\sum_{p\in\mathrm{region}}|\psi(p,n)|^{2}
 =\sum_{p\in\mathrm{region}}|\psi(p,n{+}1)|^{2}
\]
without approximation, without a Planck-scale remainder. The identity
holds for \emph{every} phase function $\varphi$ and \emph{every}
effective mass $m$ (the amplitude $|\psi|^{2}$ is mass- and
phase-blind, a consequence of the $|\psi|^{2}=\exp(-\KL)$ Gibbs
weight of Section~\ref{sec:coarsegrain}).

\subsection{Why the scope \texorpdfstring{$\HasZF$}{HasZeroFunctional}
is natural here}\label{sec:born-scope}

The $\HasZF$ hypothesis asserts that the metric-Laplacian functional
$F$ vanishes on every iterate of the Phase-1 dynamical update rule.
Under this hypothesis the induced metric is tick-invariant, hence the
KL density $\KL(g(n),g_{0},p)$ is tick-invariant, hence the Gibbs
weight $\exp(-\KL)$ is tick-invariant, hence
$|\psi(p,n)|^{2}=\exp(-\KL)$ is tick-invariant, hence the region sum
is tick-invariant. Each step is a direct rewrite.

The natural-scope diagnosis is: \textbf{Born-rule conservation and the
Schr\"odinger bound share a scope because they share a physical
input}---the tick-invariance of the KL density.
\cref{sec:T2}'s theorem concerns the \emph{form} of the per-tick
change of $\psi$; \cref{sec:born}'s concerns the \emph{conservation}
of $|\psi|^{2}$ across the tick. Both collapse trivially under
$\HasZF$ because both start from the same metric-tick-invariance.

\subsection{Theorem 3 (formal statement)}\label{sec:T3}

The headline theorem is \code{bornRuleConservation} in
\code{OmegaTheory/Emergence/BornRule.lean}, re-exported as
\code{paper\_bornRule\_conservation}:
\begin{lstlisting}
/-- **Theorem 3 (Born rule as a conservation theorem).** Under the
    static-regime hypothesis `HasZeroFunctional`, the sum
    ∑_{p ∈ region} |coarseGrainWithPhase s phase m p n|² is tick-
    invariant: probability is conserved across ticks, *as a theorem*. -/
theorem paper_bornRule_conservation
    (d : DynamicalSnapshotSequence) (hF : d.HasZeroFunctional)
    (phase : LatticePoint → ℕ → ℝ) (m : ℝ)
    (region : Finset LatticePoint) (n : ℕ) :
    (region.sum fun p =>
        (‖coarseGrainWithPhase d.toSnapshotSequence phase m p n‖) ^ 2) =
      (region.sum fun p =>
        (‖coarseGrainWithPhase d.toSnapshotSequence phase m p (n + 1)‖) ^ 2) :=
  bornRuleConservation d hF phase m region n
\end{lstlisting}

\paragraph{Faithful hypothesis inventory.}
Theorem~\ref{thm:T3-informal} shares exactly the same scoping
hypothesis as Theorem~\ref{thm:T2-informal}; the table below makes
the shared structure explicit.

\begin{longtable}{@{}p{0.32\linewidth} p{0.26\linewidth} p{0.36\linewidth}@{}}
\toprule
Hypothesis & Lean name & Meaning \\
\midrule
\endhead
$d:\DynSeq$ & --- & Phase-1 dynamical snapshot sequence. \\
$d.\HasZF$ & \code{DynamicalSnapshotSequence.HasZeroFunctional} & Shared with Theorem~\ref{thm:T2-informal}. \\
$\varphi:\mathsf{LatticePoint}\to\Nbb\to\Rbb$ & --- & Any phase function; $|\psi|^{2}$ is phase-blind, so identity holds for every $\varphi$. \\
$m:\Rbb$ & --- & Effective mass; $|\psi|^{2}$ is mass-blind. \\
$\mathrm{region}:\mathsf{Finset}\,\mathsf{LatticePoint}$ & --- & Any finite lattice region. \\
$n:\Nbb$ & --- & Tick at which we compare $n$ and $n+1$. \\
\bottomrule
\end{longtable}

\noindent Unlike Theorem~\ref{thm:T2-informal},
Theorem~\ref{thm:T3-informal} requires \textbf{no positivity
hypothesis} on the rest mass and \textbf{no Gibbs stencil
condition}---the proof uses only tick-invariance of the KL density.

\subsection{Corollaries and specialisations}\label{sec:born-coroll}

The Lean workstream exposes six additional paper-wrapper theorems:
\begin{itemize}
\item \textbf{Pointwise conservation}
  (\code{paper\_bornRule\_pointwise\_conservation}): site-wise
  identity $|\psi(p,n)|^{2}=|\psi(p,n{+}1)|^{2}$.
\item \textbf{Iterated form}
  (\code{paper\_bornRule\_amplitude\_sum\_invariant}): region sum is
  invariant across \emph{any} two ticks $n,k$.
\item \textbf{Minkowski-vacuum specialisation}
  (\code{paper\_bornRule\_on\_minkowski}): on the flat dynamical
  vacuum, conservation holds with \emph{no} side hypotheses.
\item \textbf{Honest dynamical residue bound}
  (\code{paper\_bornRule\_sum\_diff\_bound}): without any static
  hypothesis, the absolute per-tick change of the region sum of
  $|\psi|^{2}$ is bounded by the total absolute pointwise change.
\item \textbf{Residue bound in KL-density form}
  (\code{paper\_bornRule\_sum\_diff\_bound\_KL}): the same bound with
  RHS in KL Boltzmann-weight terms.
\item \textbf{Consistency corollary}
  (\code{paper\_bornRule\_sum\_eq\_of\_KL\_static}): when the KL
  density is tick-invariant on the region, the KL-form residue bound
  collapses to exact conservation.
\end{itemize}

\subsection{How Theorem~\ref{thm:T3-informal} differs from the
standard-QM Born rule}\label{sec:born-vs-standard}

Theorem~\ref{thm:T3-informal} is the \textbf{conservation} half of
Born-rule-interpreted QM, not the whole story. Standard QM takes two
independent inputs: (1) the Born postulate that $|\psi(x,t)|^{2}$ is a
probability density, and (2) unitarity of the time evolution, from
which probability conservation follows. On the V2 substrate we prove
\emph{both} at the $\HasZF$ scope: (1) is the defining identity
$|\psi|^{2}=\exp(-\KL)$ from Theorem~\ref{thm:T1}, and (2) is the
Schr\"odinger-shape bound from Theorem~\ref{thm:T2-informal}. What is
\emph{not} derived---and what belongs to a separate paper---is the
measurement-theoretic interpretation that connects $|\psi|^{2}$ to
experimental outcome frequencies (decoherence or envariance
mechanisms at the Planck scale).

\section{Non-Relativistic Limit}\label{sec:nonrel}

\subsection{From the mass shell to the kinetic term}\label{sec:nonrel-alg}

The lattice-derived mass shell $E^{2}=(pc)^{2}+(mc^{2})^{2}$ is stated
as \code{massShell\_from\_tick\_counting} in
\code{DispersionFromLattice.lean}. The standard Taylor expansion
around $p=0$
\[
E(p,m)=mc^{2}\sqrt{1+(p/mc)^{2}}
      =mc^{2}+\tfrac{p^{2}}{2m}-\tfrac{p^{4}}{8m^{3}c^{2}}+O(p^{6})
\]
identifies the first two terms as the rest energy and the classical
kinetic term. But the Lean proof does \emph{not} go through this
expansion, and the distinction matters.

\paragraph{The proof is algebraic, not analytic.}
Write
\[
A:=E(p,m)=\sqrt{(pc)^{2}+(mc^{2})^{2}},\qquad
B:=mc^{2}+\frac{p^{2}}{2m}.
\]
A one-line calculation gives the exact identity
\[
B^{2}-A^{2}=\frac{p^{4}}{4m^{2}},
\]
and $A+B\ge mc^{2}$ (from $A\ge 0$ and $B\ge mc^{2}$ when $m>0$).
Hence
\[
0\le B-A=\frac{B^{2}-A^{2}}{A+B}\le\frac{p^{4}}{4m^{2}\cdot mc^{2}}=\frac{p^{4}}{4m^{3}c^{2}}.
\]

This route avoids all epsilon-delta analytic machinery: no
derivatives, no series convergence, no limit definition enters the
proof. The Lean proof
(\code{nonRelativisticEnergy\_sq\_sub\_relativisticEnergy\_sq} followed
by \code{nonRelativistic\_energy\_approx} in
\code{DispersionBridge.lean}) reflects that structure. Compressed
into a single line: \textbf{the non-relativistic limit is an algebraic
identity, not an asymptotic expansion}.

\subsection{Theorem 4 (formal statement)}\label{sec:T4}

Proved in \code{DispersionBridge.lean} and re-exported as
\code{paper\_nonrel\_limit}:
\begin{lstlisting}
/-- **Theorem 4 (Non-relativistic limit).** For rest mass `m > 0` and
any momentum `p`, the relativistic energy differs from the non-
relativistic expression `mc² + p²/(2m)` by a quartic-in-`p` remainder:
   |relativisticEnergy p m − (m c² + p²/(2m))|
     ≤ p⁴ / (4 · m³ · c²). -/
theorem paper_nonrel_limit
    {m : ℝ} (hm : 0 < m) (p : ℝ) :
    |relativisticEnergy p m - nonRelativisticEnergy p m|
      ≤ p ^ 4 / (4 * m ^ 3 * c ^ 2) :=
  nonRelativistic_energy_approx hm p
\end{lstlisting}

Note the bound is \textbf{unconditional in $p$} (no $|p|\le\alpha mc$
hypothesis needed); the RHS grows with $p^{4}$ so at large $p$ the
bound becomes vacuous, but the \emph{inequality} is valid everywhere.
The non-relativistic \emph{regime} is the sub-region $|p|\ll mc$; the
\emph{inequality} holds globally.

The companion lemma \code{schrodinger\_is\_nonrel\_limit} provides the
formal ``Schr\"odinger $\subset$ relativistic'' hierarchy claim: if
any Schr\"odinger-side approximation produces an eigenvalue-like
quantity $K(p,m)$ with $|K-p^{2}/(2m)|\le\varepsilon$, then
$|K-(E(p,m)-mc^{2})|\le\varepsilon+p^{4}/(4m^{3}c^{2})$. Re-exported
as \code{paper\_schrodinger\_is\_nonrel\_limit}.

\subsection{Bridging Theorems~\ref{thm:T2-informal} and
\ref{thm:T3-informal}}\label{sec:nonrel-bridge}

\textbf{Composition is pointwise, not operator-level.}
Theorem~\ref{thm:T2-informal} controls the norm of the residue field
$\mathsf{schrodingerResidue}\,m\,\psi\,p\,n$ at a site $(p,n)$;
Theorem~\ref{thm:T4-informal} + \code{schrodinger\_is\_nonrel\_limit}
controls the scalar gap between an eigenvalue-like $K$ and
$p^{2}/(2m)$. Gluing the two requires identifying $K$ with the
eigenvalue of \code{discreteLaplacianC} on a \emph{plane-wave state}.

Concretely, on
$\psi_{k}:=\coarseGrainWithPhase\,s\,(\mathsf{planeWavePhase}\,k\,\omega)\,m$,
the complex discrete Laplacian is diagonal:
\[
\mathsf{discreteLaplacianC}\,\psi_{k}(p,n)=\lambda(k,\lP)\cdot\psi_{k}(p,n),
\qquad
\lambda(k,\lP)=-\frac{4}{\lP^{2}}\sum_{\mu}\sin^{2}\!\Bigl(\tfrac{k_{\mu}\,\lP}{2}\Bigr).
\]
A Taylor expansion of $\sin^{2}$ near zero gives
$|\lambda(k,\lP)-(-k^{2})|=O(\lP^{2}\,k^{4})$, which is exactly the
$\varepsilon$ that \code{schrodinger\_is\_nonrel\_limit} consumes.
Composing yields, on plane-wave test states,
\[
\bigl|\,\bigl(E(p,m)-mc^{2}\bigr)-\tfrac{p^{2}}{2m}\,\bigr|
 \le O(\lP^{2}\,k^{4})+\tfrac{p^{4}}{4m^{3}c^{2}}
\]
after identifying $p=\hbar k$.

\paragraph{What is and is not machine-checked.}
Theorems~\ref{thm:T2-informal} and \ref{thm:T4-informal} are each
machine-checked. The pointwise-to-eigenvalue reduction above is
\emph{stated} in the paper but \textbf{not yet} formalised. It is
pure lattice-Fourier analysis; the obstruction is Lean-side work, not
new mathematics. Section~\ref{sec:open} item~5 lists this honestly as
open.

\section{Two-Slit Interference}\label{sec:interference}

\subsection{What the theorem claims}\label{sec:interference-claim}

The two-slit thought experiment is the canonical signature of quantum
mechanics: two plane-wave ``paths'' coherently superpose on a common
substrate, and the resulting probability density carries a
$\cos(\Delta\varphi)$-modulated fringe pattern that is absent in
either classical path taken alone. Theorem~\ref{thm:T5-informal}
makes this signature a \emph{machine-checked theorem} of the V2
substrate: superposing two plane-wave coarse-grained fields yields,
\textbf{pointwise and without residue}, the standard interference
formula
\[
|\psi_{1}+\psi_{2}|^{2}(p,n)
 = |\psi_{1}|^{2}(p,n)+|\psi_{2}|^{2}(p,n)+2\,A(p,n)^{2}\cos\bigl(\Delta\varphi(p,n)\bigr)
\]
where $A(p,n)=\exp(-\KL(g(n),g_{0},p)/2)$ and
$\Delta\varphi(p,n)=(k_{1}-k_{2})\cdot p-(\omega_{1}-\omega_{2})\cdot n$.

\subsection{Superposition is a complex-field operation}\label{sec:interference-design}

Theorem~\ref{thm:T5-informal} deliberately lifts superposition to the
level of $\Lfield=\mathsf{LatticePoint}\to\Nbb\to\Cbb$, \textbf{not} to
the level of $\SnapSeq$. Concretely:
\begin{lstlisting}
noncomputable def superposedField
    (ψ₁ ψ₂ : LatticeComplexField) : LatticeComplexField :=
  fun p n => ψ₁ p n + ψ₂ p n
\end{lstlisting}

This is a design choice, and the honest one. \code{DiscreteMetric}
lacks a natural Lorentzian-respecting additive structure: summing two
metrics produces an object that is not, in general, a metric at all.
Any attempt to define a ``superposition of snapshot sequences'' by
lifting pointwise metric addition would introduce $O(\lP)$ residues
in the interference formula and make the two-slit identity
approximate rather than exact. The complex-field-level operation
matches the Feynman path-integral picture: two paths share
\emph{one} background and sum at the amplitude level, which is
precisely what the Lean definition captures. That the headline
identity becomes residue-free is the reward for picking the correct
venue.

\subsection{Theorem 5 (formal statement)}\label{sec:T5}

The headline theorem is \code{two\_slit\_interference} in
\code{OmegaTheory/Emergence/Interference.lean}, re-exported as
\code{paper\_two\_slit\_interference}:
\begin{lstlisting}
/-- **Theorem 5 (Two-slit interference).** For two plane-wave
    coarse-grained fields ψⱼ = coarseGrainWithPhase s
    (planeWavePhase kⱼ ωⱼ) m (j = 1, 2) on a common substrate
    s : SnapshotSequence, the superposed-field probability density is
    *exactly* |ψ₁|² + |ψ₂|² + 2·A²·cos(Δφ). -/
theorem paper_two_slit_interference
    (s : SnapshotSequence) (m : ℝ)
    (k₁ k₂ : LatticePoint) (ω₁ ω₂ : ℝ)
    (p : LatticePoint) (n : ℕ) :
    ‖superposedField (planeWaveField s m k₁ ω₁)
        (planeWaveField s m k₂ ω₂) p n‖ ^ 2 =
      ‖planeWaveField s m k₁ ω₁ p n‖ ^ 2 +
      ‖planeWaveField s m k₂ ω₂ p n‖ ^ 2 +
      2 * coarseGrainAmplitude s p n ^ 2 *
        Real.cos (planeWavePhase k₁ ω₁ p n -
                   planeWavePhase k₂ ω₂ p n) :=
  two_slit_interference s m k₁ k₂ ω₁ ω₂ p n
\end{lstlisting}

\paragraph{Faithful hypothesis inventory.}
Theorem~\ref{thm:T5-informal} is unconditional in its inputs; every
parameter is free.

\begin{longtable}{@{}p{0.32\linewidth} p{0.22\linewidth} p{0.40\linewidth}@{}}
\toprule
Hypothesis & Lean name & Meaning \\
\midrule
\endhead
$s:\SnapSeq$ & --- & Shared substrate; supplies $A(p,n)=\exp(-\KL/2)$. \\
$m:\Rbb$ & --- & Effective mass; mass-blind in $|\psi|^{2}$. \\
$k_{1},k_{2}:\mathsf{LatticePoint}$, $\omega_{1},\omega_{2}:\Rbb$ & --- & Wavevectors and frequencies of the two plane waves; fully free. \\
$p:\mathsf{LatticePoint}, n:\Nbb$ & --- & Spacetime point at which the identity is evaluated. \\
\bottomrule
\end{longtable}

\noindent No positivity of mass, no Gibbs stencil condition, no
$\HasZF$. The identity holds in the free algebraic regime.

\subsection{Corollaries}\label{sec:T5-coroll}

Two subsequent algebraic rewrites deliver the textbook fringe shape:
\begin{itemize}
\item \textbf{Boltzmann-weight form}
  (\code{paper\_two\_slit\_interference\_KL}): rewriting $|\psi_{j}|^{2}$
  via \code{planeWaveField\_abs\_sq} gives
  $|\psi_{1}+\psi_{2}|^{2}(p,n)=2\exp(-\KL)+2A^{2}\cos(\Delta\varphi)$.
\item \textbf{Canonical visibility} (\code{paper\_double\_slit\_visibility}):
  because both plane-wave fields inherit the \emph{same} substrate
  amplitude $A$ by construction,
  $|\psi_{1}+\psi_{2}|^{2}(p,n)=2A^{2}(1+\cos\Delta\varphi)=4A^{2}\cos^{2}(\Delta\varphi/2)$.
\item \textbf{Constructive peak}
  (\code{paper\_two\_slit\_constructive\_peak}): at $\Delta\varphi=0$
  the intensity is $4A^{2}$, four times the single-slit intensity.
\item \textbf{Destructive null}
  (\code{paper\_two\_slit\_destructive\_null}): at $\Delta\varphi=\pi$
  the intensity is identically zero---exact cancellation.
\item \textbf{Flat-vacuum sanity} (\code{paper\_two\_slit\_on\_flat}):
  on the Minkowski vacuum where $A=\exp(-1)$ pointwise, the pattern
  collapses to $2\exp(-2)(1+\cos\Delta\varphi)$.
\end{itemize}

\subsection{Compatibility with Born-rule conservation}\label{sec:T5-born-compat}

Under $\HasZF$, the substrate amplitude $A(p,n)$ is
\emph{tick-invariant}: Theorem~\ref{thm:T3-informal} implies the KL
density is constant in $n$, hence so is $A$. The formal consequence
for interference is \code{paper\_two\_slit\_envelope\_tick\_invariant}:
\begin{lstlisting}
theorem paper_two_slit_envelope_tick_invariant
    (d : DynamicalSnapshotSequence) (hF : d.HasZeroFunctional)
    (m : ℝ) (k₁ k₂ : LatticePoint) (ω₁ ω₂ : ℝ)
    (p : LatticePoint) (n k : ℕ) :
    coarseGrainAmplitude d.toSnapshotSequence p n =
      coarseGrainAmplitude d.toSnapshotSequence p k
\end{lstlisting}

Physically: \textbf{probability is conserved at the substrate
(amplitude) level, while interference fringes still move across ticks}---
because the phase difference
$\Delta\varphi(p,n)=(k_{1}-k_{2})\cdot p-(\omega_{1}-\omega_{2})\cdot n$
picks up $(\omega_{2}-\omega_{1})\cdot n$ per tick.

\subsection{What Theorem~\ref{thm:T5-informal} does and does not
buy}\label{sec:T5-neg}

It \emph{is} a machine-checked derivation of the textbook two-slit
interference identity from a discrete gravitational substrate
(best-effort framing). It is \emph{not} a derivation of the geometric
origin of the plane-wave phase---the identity takes $(k_{j},\omega_{j})$
as external inputs via \code{planeWavePhase}. A principled derivation
of the outgoing $k$ from the geometry of the ``slit'' pattern on the
lattice is a separate workstream (Section~\ref{sec:open}, item~2).

\section{Heisenberg Uncertainty}\label{sec:heisenberg}

\subsection{The claim}\label{sec:heisenberg-claim}

The Heisenberg uncertainty principle
$\Delta x\cdot\Delta p\ge\hbar/2$ is, in standard non-relativistic QM,
a consequence of the Robertson inequality~\cite{Robertson1929} applied
to the canonical commutation relation
$[\xhat,\phat]=i\hbar$~\cite{Heisenberg1927}.
Theorem~\ref{thm:T6-informal} derives the Heisenberg inequality on
the V2 substrate: position $\xhat_{\mu}$ and momentum
$\phat_{\mu}=-i\hbar\partial_{\mu}$ are given concrete definitions on
the lattice, the Robertson inequality holds as a \emph{pure
Cauchy--Schwarz fact} on any $\Lfield$, and the canonical
$(\hbar/2)^{2}$ lower bound follows when the lattice commutator
$[\xhat_{\mu},\phat_{\mu}]$ evaluated against the state's own
expectation values matches its continuum value $\hbar$.

\subsection{Lattice operators and their
variances}\label{sec:heisenberg-ops}

\code{Heisenberg.lean} defines a \code{LatticeOperator} as a map
$\Lfield\to\mathsf{LatticePoint}\to\Nbb\to\Cbb$---a pointwise complex
transformation of a wavefunction. The position and momentum operators
are:
\begin{lstlisting}
noncomputable def positionOperator (μ : Fin 4) : LatticeOperator :=
  fun ψ p n => ((p μ : ℝ) : ℂ) * ψ p n

noncomputable def momentumOperator (μ : Fin 4) : LatticeOperator :=
  fun ψ p n => -Complex.I * ((hbar : ℝ) : ℂ) *
    ((ψ (shiftFin p μ) n - ψ (shiftBackFin p μ) n) / ((2 * l_P : ℝ) : ℂ))
\end{lstlisting}

The position operator multiplies the wavefunction by the $\mu$-th
integer coordinate of $p$, reflecting the lattice embedding
$\xhat_{\mu}\psi=p_{\mu}\cdot\psi$. The momentum operator is the
symmetric-difference approximation to $-i\hbar\,\partial/\partial x_{\mu}$,
using forward / backward lattice shifts and dividing by $2\lP$.

The variance definitions:
\begin{lstlisting}
noncomputable def variance_x
    (μ : Fin 4) (ψ : LatticeComplexField)
    (region : Finset LatticePoint) (tick : ℕ) : ℝ :=
  l2Variance (positionOperator μ) (expValue (positionOperator μ) ψ region tick)
    ψ region tick

noncomputable def variance_p
    (μ : Fin 4) (ψ : LatticeComplexField)
    (region : Finset LatticePoint) (tick : ℕ) : ℝ :=
  l2Variance (momentumOperator μ) (expValue (momentumOperator μ) ψ region tick)
    ψ region tick
\end{lstlisting}

Both are non-negative
(\code{variance\_x\_nonneg}, \code{variance\_p\_nonneg}). The
generalised $\mathsf{l2Variance}\,A\,\mu\,\psi\,R\,t$ is
$\sum_{p\in R}\|A\psi(p,t)-\mu\psi(p,t)\|^{2}$---the $L^{2}$ distance
of $A\psi$ from its candidate mean $\mu\psi$. When $\mu=\langle A\rangle$
this recovers the textbook variance; the generalisation is needed for
the Robertson Cauchy--Schwarz step.

\subsection{The
\texorpdfstring{$\CMM$}{CommutatorMatchesMean} hypothesis}\label{sec:CMM}

The Heisenberg bound is proved conditional on a hypothesis that
encodes the continuum commutator value on the lattice:
\begin{lstlisting}
def CommutatorMatchesMean
    (A B : LatticeOperator) (c : ℝ) (ψ : LatticeComplexField)
    (region : Finset LatticePoint) (tick : ℕ) : Prop :=
  2 * (robertsonCrossTerm A B
        (expValue A ψ region tick)
        (expValue B ψ region tick) ψ region tick).im = c
\end{lstlisting}

Physically, this says that
$2\cdot\mathrm{Im}\langle(A-\langle A\rangle)\psi\mid(B-\langle B\rangle)\psi\rangle=c$:
the expectation value of $-i[A,B]$ on $\psi$, evaluated at the
variance-centred operators, equals $c$. Setting $A=\xhat_{\mu}$,
$B=\phat_{\mu}$, $c=\hbar$ then says the lattice-evaluated commutator
reproduces the continuum value $\hbar$ when integrated against $\psi$
across the region.

\paragraph{Scope of the hypothesis.}
On a smooth state (plane wave, Gaussian) sampled at lattice sites,
$\CMM$ holds exactly in the $\lP\to 0$ limit. On the lattice at
finite $\lP$, the symmetric-difference momentum operator satisfies
$[\xhat,\phat]\psi(p)=i\hbar\cdot(\psi(p+\hat e)+\psi(p-\hat e))/2$
rather than $i\hbar\cdot\psi(p)$, so the commutator picks up an
$O(\lP^{2})$ correction for states smooth on the lattice scale. The
$\CMM$ predicate is the honest \emph{scope of the closure}. A
quantitative \code{lattice\_commutator\_converges} theorem bounding
the $O(\lP^{2})$ slack is future work.

\subsection{Theorem 6 (formal statement)}\label{sec:T6}

\begin{lstlisting}
/-- **Theorem 6 (Heisenberg uncertainty, from the lattice).** For
    any direction μ, any ψ, any finite region R, tick n, and
    normalised state (∑ |ψ|² = 1) satisfying `CommutatorMatchesMean
    (x̂_μ) (p̂_μ) ℏ ψ R n`, the variance product is bounded below
    by `(ℏ/2)²`. -/
theorem paper_heisenberg_uncertainty
    (μ : Fin 4) (ψ : LatticeComplexField)
    (region : Finset LatticePoint) (tick : ℕ)
    (hnorm : ∑ p ∈ region, ‖ψ p tick‖ ^ 2 = 1)
    (hcomm : CommutatorMatchesMean (positionOperator μ)
              (momentumOperator μ) hbar ψ region tick) :
    variance_x μ ψ region tick * variance_p μ ψ region tick ≥
      (hbar / 2) ^ 2
\end{lstlisting}

\paragraph{Faithful hypothesis inventory.}

\begin{longtable}{@{}p{0.22\linewidth} p{0.26\linewidth} p{0.42\linewidth}@{}}
\toprule
Hypothesis & Lean name & Meaning \\
\midrule
\endhead
$\mu:\mathsf{Fin}\,4$ & --- & Direction index. \\
$\psi:\Lfield$ & --- & The state under consideration. \\
$\mathrm{region}:\mathsf{Finset}\,\mathsf{LatticePoint}$ & --- & Finite lattice region. \\
$\mathrm{tick}:\Nbb$ & --- & Time slice. \\
\code{hnorm} & --- & $L^{2}$-normalisation on the region. Not used in the proof (Robertson is scale-free) but required for the standard physical reading. \\
\code{hcomm} & $\CMM$ & Honest lattice-continuum commutator-matching hypothesis. \\
$0<\hbar$ & \code{hbar\_pos} & Positivity of $\hbar$ (Planck-constant axiom, not user-side). \\
\bottomrule
\end{longtable}

\noindent The conclusion
$\mathrm{variance}_{x}\cdot\mathrm{variance}_{p}\ge(\hbar/2)^{2}$ is
exact on-regime: under $\CMM(\cdot)(\cdot)\,\hbar$, the bound has the
classic textbook form with no Planck-scale remainder at the level of
the inequality itself; the remainder lives inside the \code{hcomm}
hypothesis, as designed.

\subsection{Corollaries}\label{sec:T6-coroll}

\begin{itemize}
\item \textbf{Abstract Robertson form}
  (\code{paper\_heisenberg\_uncertainty\_abstract}): for any two
  lattice operators $A,B$ with $\mathsf{CommutatorMatchesValue}\,A\,B\,c$,
  $\mathsf{l2Variance}\,A\,0\,\psi\,R\,n\cdot\mathsf{l2Variance}\,B\,0\,\psi\,R\,n\ge(c/2)^{2}$.
\item \textbf{Positive-variance corollaries}
  (\code{variance\_x\_pos\_of\_heisenberg},
  \code{variance\_p\_pos\_of\_heisenberg}). Under $\CMM\cdot\hbar$,
  both variances must be strictly positive---``no simultaneously
  localised position and momentum''.
\item \textbf{Robertson exact form} (\code{robertson\_uncertainty\_exact}):
  without the commutator-matching hypothesis,
  $\mathsf{l2Variance}\,A\,\mu_{A}\cdot\mathsf{l2Variance}\,B\,\mu_{B}\ge|\mathrm{Im}\,\langle(A-\mu_{A})\psi\mid(B-\mu_{B})\psi\rangle|^{2}$
  for \emph{any} lattice operators and means.
\end{itemize}

\subsection{What Theorem~\ref{thm:T6-informal} does and does not
say}\label{sec:T6-neg}

It \emph{is} a machine-checked derivation of the Heisenberg bound
from a discrete gravitational substrate, under an honestly-scoped
commutator-matching hypothesis (best-effort framing). It is
\emph{not} a derivation that makes the lattice commutator equal its
continuum value as a theorem---that would require a quantitative
$O(\lP^{2})$ Taylor bound on smooth test states, which is not yet
formalised (Section~\ref{sec:open}, item~4). It is \emph{not} an
uncertainty relation for continuous Schr\"odinger operators on
$L^{2}(\Rbb^{4})$---the theorem is stated inside the lattice
framework.

\section{Measurement / Collapse}\label{sec:measurement}

\subsection{The claim}\label{sec:measurement-claim}

In standard non-relativistic QM, the measurement / wavefunction-
collapse postulate is a \emph{separate axiom} from the Schr\"odinger
evolution: upon measurement of an observable with a discrete
eigenvalue basis, the wavefunction is projected onto the eigenstate
of the measured outcome and renormalised to unit norm, with the
probability of outcome $q$ given by the Born ratio $|\psi(q)|^{2}$.
The four-clause package---probability is Born-ratio, unit norm
post-measurement, support concentration at $q$, unit modulus at
$q$---constitutes the collapse postulate.

Theorem~\ref{thm:T7-informal} delivers this package as a
\emph{theorem} on the V2 substrate. The projector is a pointwise
indicator (\code{measurementProjector}), the renormalisation is
division by a region-relative $L^{2}$-norm
(\code{normaliseCoarseGrain}), and the composition
\code{postMeasurementState\,region\,$\psi$\,q\,tick} satisfies all
four clauses of the textbook postulate, with the crucial honest
acknowledgement that the post-measurement state falls outside the
image of \code{coarseGrain}---collapse is non-unitary, and this is a
theorem rather than an axiom.

\subsection{Region-relative \texorpdfstring{$L^{2}$}{L2}
norm}\label{sec:measurement-norm}

Because the Omega substrate lives on $\Zbb^{4}$ with no natural total
$\ell^{2}$ mass, every measurement-postulate statement is parametrised
by a $\mathsf{Finset}\,\mathsf{LatticePoint}$ region on which the
wavefunction is conditioned. The region-relative norm is:
\begin{lstlisting}
noncomputable def regionL2NormSq
    (region : Finset LatticePoint) (ψ : LatticeComplexField)
    (tick : ℕ) : ℝ :=
  ∑ p ∈ region, ‖ψ p tick‖ ^ 2
\end{lstlisting}

with $\mathsf{regionL2Norm}:=\sqrt{\mathsf{regionL2NormSq}}$. This is
the standard Born-probability denominator.

\subsection{The measurement operation}\label{sec:measurement-op}

\code{Measurement.lean} composes three pointwise operations:
\begin{itemize}
\item \textbf{Projector}
  ($\mathsf{measurementProjector}\,q\,\psi\,p\,n:=\mathbf{1}[p=q]\cdot\psi(p,n)$).
\item \textbf{Renormaliser}
  ($\mathsf{normaliseCoarseGrain}\,R\,\psi\,t\,p\,n:=\psi(p,n)/\mathsf{regionL2Norm}\,R\,\psi\,t$).
\item \textbf{Composition} (\code{postMeasurementState}).
\end{itemize}

\subsection{Theorem 7 (four-clause bundle)}\label{sec:T7}

\begin{lstlisting}
/-- **Theorem 7 (Measurement postulate).** Let region, ψ, q, tick
    be as above with q ∈ region and ψ q tick ≠ 0. Then:
    (1) probability_of_outcome = Born ratio,
    (2) unit norm post-measurement,
    (3) support concentrated at q,
    (4) unit modulus at q. -/
theorem paper_measurement_postulate
    (region : Finset LatticePoint) (ψ : LatticeComplexField)
    (q : LatticePoint) (tick : ℕ)
    (hq : q ∈ region) (hψq : ψ q tick ≠ 0) :
    probability_of_outcome region ψ q tick =
        (‖ψ q tick‖) ^ 2 / regionL2NormSq region ψ tick
    ∧ regionL2NormSq region (postMeasurementState region ψ q tick) tick = 1
    ∧ (∀ p' : LatticePoint, p' ≠ q →
        postMeasurementState region ψ q tick p' tick = 0)
    ∧ ‖postMeasurementState region ψ q tick q tick‖ = 1
\end{lstlisting}

\paragraph{Faithful hypothesis inventory.}

\begin{longtable}{@{}p{0.28\linewidth} p{0.22\linewidth} p{0.44\linewidth}@{}}
\toprule
Hypothesis & Lean name & Meaning \\
\midrule
\endhead
$\mathrm{region}$ & --- & Finite region on which probability is conditioned. \\
$\psi$ & --- & Pre-measurement state. \\
$q$ & --- & Outcome of the measurement. \\
$\mathrm{tick}$ & --- & Time slice at which the measurement occurs. \\
$q\in\mathrm{region}$ & \code{hq} & Outcome is inside conditioned region. \\
$\psi(q,\mathrm{tick})\ne 0$ & \code{h$\psi$q} & Non-zero pre-measurement amplitude; otherwise the Born ratio is undefined. \\
\bottomrule
\end{longtable}

\noindent The theorem has no side hypothesis on the substrate (no
$\HasZF$, no Gibbs stencil, no mass positivity). It holds for
\emph{any} $\Lfield$. This matches the physics: the collapse
postulate is a statement about what happens \emph{upon} measurement;
it is independent of how the wavefunction got there.

\subsection{Non-unitarity: the honest
acknowledgement}\label{sec:measurement-nonunitary}

A foundational question forced by Theorem~\ref{thm:T7-informal} is:
does the post-measurement state admit a Schr\"odinger-evolution
treatment? The honest answer is \textbf{no}, and the paper commits to
this openly.

\begin{lstlisting}
/-- **Collapse is non-unitary.** For any p' ≠ q,
    postMeasurementState region ψ q tick p' tick ≠ coarseGrain s p' tick
    for any SnapshotSequence s. -/
theorem paper_postMeasurement_non_unitary
    (region : Finset LatticePoint) (ψ : LatticeComplexField)
    (q : LatticePoint) (tick : ℕ)
    (p' : LatticePoint) (hp'_ne : p' ≠ q)
    (s : SnapshotSequence) :
    postMeasurementState region ψ q tick p' tick ≠
      coarseGrain s p' tick
\end{lstlisting}

Proof: the post-measurement state is zero at $p'$ by support
concentration (clause~3), while \code{coarseGrain\,s\,p'\,tick} has
strictly positive modulus because $A=\exp(-\KL/2)$ is strictly
positive everywhere. Hence they disagree at $p'$.

\paragraph{Physical reading.}
Theorem~\ref{thm:T2-informal} is a theorem about a specific
well-behaved class of $\Lfield$s---those obtained by coarse-graining
a snapshot sequence. Measurement takes us out of this class by
design: \code{coarseGrain\,s} has everywhere-positive modulus, while
the collapsed state has zero modulus off $q$.
Theorem~\ref{thm:T2-informal} therefore \emph{does not apply} to the
post-measurement state, and it \emph{should not}---wavefunction
collapse is non-unitary, and the substrate formalism recognises this.

\subsection{Compatibility with Born-rule
conservation}\label{sec:measurement-compat}

Under $\HasZF$, the Born-ratio probability is itself tick-invariant:
\begin{lstlisting}
theorem paper_probability_of_outcome_invariant_under_hasZero
    (d : DynamicalSnapshotSequence) (hF : d.HasZeroFunctional)
    (phase : LatticePoint → ℕ → ℝ) (m : ℝ)
    (region : Finset LatticePoint) (q : LatticePoint) (n k : ℕ) :
    probability_of_outcome region
        (coarseGrainWithPhase d.toSnapshotSequence phase m) q n =
      probability_of_outcome region
        (coarseGrainWithPhase d.toSnapshotSequence phase m) q k
\end{lstlisting}

The Minkowski-vacuum specialisation
\code{paper\_probability\_of\_outcome\_flat} says every outcome in the
region is uniformly distributed:
$\mathsf{probability\_of\_outcome}=1/|\mathrm{region}|$.

\subsection{What Theorem~\ref{thm:T7-informal} does and does not
buy}\label{sec:T7-neg}

It \emph{is} a machine-checked derivation of the four-clause collapse
postulate (projection, renormalisation, Born ratio, support
concentration), with the non-unitarity of collapse as a theorem
rather than a side remark (best-effort framing). It is \emph{not} a
resolution of the measurement problem. It derives the
\emph{phenomenology} of collapse but not \emph{why} measurements
happen or \emph{how} a specific outcome is selected. That deeper
layer---continuous-measurement decoherence / envariance---remains
beyond the scope of this paper.

\section{Entanglement and Bell Inequality
Violation}\label{sec:entanglement}

\subsection{The dramatic claim}\label{sec:entanglement-claim}

Bell's theorem~\cite{Bell1964,CHSH1969} states that no local hidden
variable (LHV) theory can reproduce the correlations of certain
entangled quantum states---the classical Bell bound
$|S_{\mathrm{CHSH}}|\le 2$ is provably violated by suitable choices of
entangled state and measurement angles, with the quantum Tsirelson
bound $|S_{\mathrm{CHSH}}|\le 2\sqrt{2}$~\cite{Tsirelson1980} as the
ultimate upper limit. This is the mathematical foundation of
Einstein's ``spooky action at a distance'' and the experimental
closure of LHV theories~\cite{Aspect1982,Hensen2015,Giustina2015}.

Theorem~\ref{thm:T8-informal} reconstructs this result on the V2
substrate, with four machine-checked conclusions:
\begin{enumerate}
\item The Bell state $\BellField=(|00\rangle+|11\rangle)/\sqrt{2}$---
  built as an explicit
  $\mathsf{TwoBodyField}:\mathsf{LatticePoint}\to\mathsf{LatticePoint}\to\Nbb\to\Cbb$
  on the lattice---is \textbf{entangled}, not the tensor product of
  any two single-body fields (\code{bellField\_isEntangled}).
\item Its two-particle correlator is \textbf{exactly}
  $\cos(\alpha-\beta)$---not merely bounded by 1
  (\code{cos\_correlation\_theorem}).
\item At the canonical CHSH angles $(0,\pi/2,\pi/4,3\pi/4)$, the CHSH
  value \textbf{attains the Tsirelson bound} $2\sqrt{2}$
  (\code{chsh\_tsirelson\_bell}).
\item The same CHSH value \textbf{strictly exceeds} the classical
  Bell bound $2$ (\code{bell\_inequality\_violation}).
\end{enumerate}

Together these four statements constitute a machine-checked
formalisation of Bell's theorem on the discrete-gravity lattice.

\subsection{Two-body fields and tensor
products}\label{sec:entanglement-2body}

\code{Entanglement.lean} introduces the two-body field type and
tensor-product operation:
\begin{lstlisting}
noncomputable def tensorProduct
    (ψA ψB : LatticeComplexField) : TwoBodyField :=
  fun p q n => ψA p n * ψB q n

notation:70 ψA " ⊗ₜ " ψB => tensorProduct ψA ψB
\end{lstlisting}

The tensor product is bilinear (\code{tensorProduct\_add\_left},
\code{tensorProduct\_add\_right}) and factorises the squared modulus
(\code{tensorProduct\_abs\_sq}). The factorisability predicate:
\begin{lstlisting}
def IsProduct (Ψ : TwoBodyField) : Prop :=
  ∃ ψA ψB : LatticeComplexField, Ψ = ψA ⊗ₜ ψB

def IsEntangled (Ψ : TwoBodyField) : Prop :=
  ¬ IsProduct Ψ
\end{lstlisting}

\paragraph{Design note.}
As with two-slit interference
(Section~\ref{sec:interference-design}), superposition here lives at
the complex-field level (\code{TwoBodyField} is a pointwise
complex-valued function), \textbf{not} at the snapshot-sequence /
metric level. The Bell theorem's residue-free form is a reward for
the honest venue choice.

\subsection{The Bell state on the lattice}\label{sec:bell-lattice}

The Bell state $|\Phi^{+}\rangle=(|00\rangle+|11\rangle)/\sqrt{2}$ is
realised on the lattice by picking two distinguished lattice points
$\mathsf{zeroPoint}$ and $\mathsf{onePoint}$ and defining:
\begin{lstlisting}
noncomputable def bellField : TwoBodyField :=
  fun p q n =>
    if p = zeroPoint ∧ q = zeroPoint then (1 / Real.sqrt 2 : ℂ)
    else if p = onePoint ∧ q = onePoint then (1 / Real.sqrt 2 : ℂ)
    else 0
\end{lstlisting}

The four diagonal / off-diagonal values confirm amplitude $1/\sqrt{2}$
on $(0,0)$ and $(1,1)$, zero on $(0,1),(1,0)$, zero everywhere else.

\subsection{Theorem 8a --- structural entanglement}\label{sec:T8a}

\begin{lstlisting}
/-- **Theorem 8a (Bell state is entangled).** There exist no
    ψA, ψB such that bellField = ψA ⊗ ψB. -/
theorem paper_bell_state_exists :
    IsEntangled bellField :=
  bellField_isEntangled
\end{lstlisting}

Proof sketch: if $\BellField=\psi_{A}\otimes\psi_{B}$, then
evaluating at $(0,0)$ gives both factors non-zero at $0$; evaluating
at $(0,1)$ forces $\psi_{B}(1)=0$; evaluating at $(1,1)$ requires
$\psi_{A}(1)\cdot\psi_{B}(1)=1/\sqrt{2}$---contradiction.

\subsection{Theorem 8b --- the cosine
correlator}\label{sec:T8b}

\begin{lstlisting}
/-- **Theorem 8b (cos correlation).** The Bell-state correlator
    is exactly cos(α − β), not merely bounded by 1. -/
theorem paper_cos_correlation (α β : ℝ) :
    correlationBell α β = Real.cos (α - β) :=
  cos_correlation_theorem α β
\end{lstlisting}

Standard consistency checks: aligned-angle perfect correlation,
antipodal-angle perfect anti-correlation, symmetry, $|E|\le 1$.

\subsection{Theorem 8c --- the Tsirelson bound
attained}\label{sec:T8c}

Define the CHSH-Tsirelson functional:
\begin{lstlisting}
noncomputable def chshTsirelson (E : ℝ → ℝ → ℝ)
    (a a' b b' : ℝ) : ℝ :=
  E a b - E a b' + E a' b + E a' b'

noncomputable def chshTsirelsonBell (a a' b b' : ℝ) : ℝ :=
  chshTsirelson correlationBell a a' b b'
\end{lstlisting}

At the canonical Bell-state-optimal angles $(0,\pi/2,\pi/4,3\pi/4)$:
\begin{lstlisting}
/-- **Theorem 8c (Tsirelson bound attained).** At canonical
    angles, CHSH = 2·√2 exactly. -/
theorem paper_chsh_tsirelson_bound :
    chshTsirelsonBell 0 (Real.pi / 2) (Real.pi / 4) (3 * Real.pi / 4) =
      2 * Real.sqrt 2 :=
  chsh_tsirelson_bell
\end{lstlisting}

The proof expands the four correlators using
\code{cos\_correlation\_theorem}:
$E(0,\pi/4)=\sqrt{2}/2$, $E(0,3\pi/4)=-\sqrt{2}/2$,
$E(\pi/2,\pi/4)=\sqrt{2}/2$, $E(\pi/2,3\pi/4)=\sqrt{2}/2$. Then
$\mathsf{chshTsirelson}$ gives
$\sqrt{2}/2-(-\sqrt{2}/2)+\sqrt{2}/2+\sqrt{2}/2=2\sqrt{2}$.

\paragraph{Physical interpretation.}
$2\sqrt{2}$ is the \emph{Tsirelson bound}: the maximum CHSH value
attainable by \emph{any} quantum correlator. $\BellField$ saturates
it at the canonical angles---the Bell state is the extremal quantum
correlator.

\subsection{Theorem 8d --- classical Bell bound
violated}\label{sec:T8d}

\begin{lstlisting}
/-- **Theorem 8d (Bell inequality violation).** CHSH > 2. -/
theorem paper_bell_inequality_violation :
    chshTsirelsonBell 0 (Real.pi / 2) (Real.pi / 4) (3 * Real.pi / 4) > 2 :=
  bell_inequality_violation
\end{lstlisting}

Proof: $2\sqrt{2}>2$ because $\sqrt{2}>1$. The classical Bell bound
$|S_{\mathrm{CHSH}}|\le 2$ is \emph{strictly} exceeded by
$\BellField$. No local hidden variable theory can reproduce these
correlations.

\subsection{Faithful hypothesis inventory}\label{sec:T8-hyp}

Theorems~8a--8d are all \textbf{unconditional}: no side hypotheses
beyond the definitions of $\BellField$,
$\mathsf{correlationBell}$, $\mathsf{chshTsirelson}$.

\begin{longtable}{@{}p{0.16\linewidth} p{0.48\linewidth} p{0.30\linewidth}@{}}
\toprule
Theorem & Lean name & Hypotheses \\
\midrule
\endhead
8a & \code{paper\_bell\_state\_exists} & none \\
8b & \code{paper\_cos\_correlation} & none beyond free parameters $\alpha,\beta$ \\
8c & \code{paper\_chsh\_tsirelson\_bound} & none (angles are explicit constants) \\
8d & \code{paper\_bell\_inequality\_violation} & none \\
\bottomrule
\end{longtable}

\subsection{What Theorem~\ref{thm:T8-informal} does and does not
say}\label{sec:T8-neg}

It \emph{is} a machine-checked derivation of a Tsirelson-bound-
attaining entangled state on a discrete gravitational substrate
(best-effort framing). It is \emph{not} a derivation of the Born
rule on two-body states (that identification is implicit in
$\mathsf{correlationBell}\,\alpha\,\beta=\cos(\alpha-\beta)$ as the
expected product of $\pm 1$ outcomes). It is \emph{not} an
experimental-loophole-free claim (that is an empirical question). It
is \emph{not} an independent derivation of special relativity or
quantum field theory. The relativistic bridge is the subject of
\cref{sec:kleingordon}.

\section{Relativistic Extension: Klein--Gordon Residue
  (Theorem~9)}\label{sec:kleingordon}

\subsection{What the theorem claims}

Sections \ref{sec:schrodinger}--\ref{sec:bell} derive the eight
non-relativistic QM postulates of the capstone record from the V2
substrate. Theorem~9 takes the same coarse-graining map $\Lmap$ and
the same dynamical snapshot sequence and asks the relativistic
question: does the coarse-grained field satisfy the Klein--Gordon
equation $\bigl((1/c^{2})\partial_{t}^{2} -
\Delta_{\mathrm{lat}}\bigr)\psi + (m^{2}c^{2}/\hbar^{2})\psi = 0$ on
the lattice?

The honest answer is: \emph{up to an explicit, Planck-scale, UV cutoff
error}. Concretely, define
\begin{align*}
\mathrm{KG}[\psi](p, n) \;\coloneqq\;&\;
  \tfrac{1}{c^{2}}\,\partial_{t}^{2,\mathrm{disc}}\psi(p, n)
  - \Delta_{\mathrm{lat}}\psi(p, n) \\
&\;+ \tfrac{m^{2} c^{2}}{\hbar^{2}}\,\psi(p, n+1),
\end{align*}
where $\partial_{t}^{2,\mathrm{disc}}\psi(p, n) \coloneqq
\bigl(\psi(p, n+2) - 2\psi(p, n+1) + \psi(p, n)\bigr)/\tP^{2}$ is the
centred second difference over ticks $n, n+1, n+2$ normalised by
$\tP^{2}$, and $\Delta_{\mathrm{lat}}$ is the five-point complex
Laplacian from \cref{sec:schrodinger}. The theorem is:

\begin{theorem}[Relativistic KG residue bound;
  \code{coarseGrain\_satisfies\_kleinGordon\_dynamic}]
\label{thm:T9-informal}
On a $\DynSeq$ $d$ with $d.\HasZF$, a Gibbs background, for every rest
mass $m \in \Rbb$, for every lattice point $p$, and for every interior
tick $n+1$,
\[
\bigl\|\mathrm{KG}[\coarseGrain\,d.\mathsf{toSnapshotSequence}](p, n)\bigr\|
 \;\le\; \frac{16}{\lP^{2}} + \frac{m^{2} c^{2}}{\hbar^{2}},
\]
with bound constant
$\code{kleinGordonBoundConst}\;m = 16/\lP^{2} + m^{2} c^{2}/\hbar^{2}$.
\end{theorem}

\subsection{The UV cutoff is a feature of the lattice, not a defect of
  the proof}

Unlike the Schr\"odinger residue of \cref{thm:T2-informal} (which is
$\mathcal{O}(\lP)$ and vanishes in the continuum limit), the
Klein--Gordon residue bound is \emph{not} Planck-small. The leading
term $16/\lP^{2}$ is the Planck-scale UV cutoff of the discrete
Laplacian on a Planck-scale lattice: $\Delta_{\mathrm{lat}}$ is
genuinely of order $1/\lP^{2}$ as an operator norm, and any
machine-checked residue bound for a second-order differential operator
on a lattice with spacing $\lP$ inherits that order of magnitude.

This is not a defect of the proof, and we emphasise the framing:

\begin{itemize}[leftmargin=1.4em]
\item The continuum relation $(\Box + m^{2} c^{2}/\hbar^{2})\psi = 0$
is the $\lP \to 0$ limit of the discrete equation, not the discrete
equation itself. Recovering zero residue requires a
$C^{4}$-bounded coarse-graining over many lattice sites, which is
deferred to a separate workstream
(\cref{sec:open} item~\ref{open:kg-coarse}).

\item The $16/\lP^{2}$ constant is an \emph{explicit prediction} of
Omega-theory's discrete-spacetime ontology: the UV cutoff is not a
regulator to be removed, it is a physical scale set by the Planck
length. QFT on a smooth Minkowski background has no analogue of this
constant; that is precisely the operational signature by which
Omega-theory differs from pure relativistic QFT.

\item The mass contribution $m^{2} c^{2}/\hbar^{2}$ carries the
continuum Klein--Gordon mass term verbatim. Since the substrate
supplies $\hbar = \EP \cdot \tP$ from
\code{Spacetime/Constants.lean}, the mass term is \emph{fixed} by the
substrate, not inserted.
\end{itemize}

\subsection{Theorem 9 (formal statement)}

Proved in \code{OmegaTheory/Emergence/KleinGordon.lean} as
\code{coarseGrain\_satisfies\_kleinGordon\_dynamic}:

\begin{lstlisting}
/-- **Theorem 9 (Klein--Gordon residue bound, dynamical).**
    On a DynamicalSnapshotSequence `d` with `d.HasZeroFunctional`
    and a Gibbs stencil at tick `n+1`, for every rest mass `m : ℝ`
    and every lattice point `p`, the Klein--Gordon residue of the
    coarse-grained field is bounded by
    `kleinGordonBoundConst m = 16/ℓ_P² + m²c²/ℏ²`. -/
theorem coarseGrain_satisfies_kleinGordon_dynamic
    (d : DynamicalSnapshotSequence) (hF : d.HasZeroFunctional)
    (p : LatticePoint) (n : ℕ) (m : ℝ)
    (hcenter : 0 ≤ informationDensityKL
                    (d.toSnapshotSequence.metric (n + 1))
                    d.toSnapshotSequence.reference p)
    (hstencil : ∀ μ : Fin 4,
      0 ≤ informationDensityKL
            (d.toSnapshotSequence.metric (n + 1))
            d.toSnapshotSequence.reference (shiftFin p μ) ∧
      0 ≤ informationDensityKL
            (d.toSnapshotSequence.metric (n + 1))
            d.toSnapshotSequence.reference (shiftBackFin p μ)) :
    ‖kleinGordonResidue m (coarseGrain d.toSnapshotSequence) p n‖
      ≤ kleinGordonBoundConst m
\end{lstlisting}

\paragraph{Faithful hypothesis inventory.}

\begin{table}[h]
\small
\begin{tabular}{@{}l l p{0.46\linewidth}@{}}
\toprule
Hypothesis & Lean name & Meaning \\
\midrule
$d : \DynSeq$ & --- & Phase-1 dynamical snapshot sequence. \\
$d.\HasZF$ & \code{DynamicalSnapshotSequence.HasZero-}\code{Functional} & Shared scope with T2 and T3. Under this hypothesis Phase~1's \code{static\_reduces\_to\_snapshot\_sequence} collapses the metric to tick-invariant; the KG second time-difference then vanishes by \code{secondTickDiffC\_coarseGrain\_static}, and no multi-tick Gibbs hypothesis is needed. \\
$m : \Rbb$ & --- & Rest mass; unlike T2, \emph{no} positivity hypothesis on $m$ is required (the KG equation handles $m = 0$ as the lightcone case). \\
$0 \le \IKL(g(n{+}1), g_{0}, p)$ & \code{hcenter} & Pointwise Gibbs at the central lattice point at tick $n{+}1$. \\
$\forall \mu.\; 0 \le \IKL(g(n{+}1), g_{0}, p \pm e_{\mu})$ & \code{hstencil} & Gibbs on the $8$ nearest neighbours at tick $n{+}1$. \textbf{Same 9-point spatial stencil as T2} --- the $\HasZF$ scope collapses any multi-tick time-stencil to this single-tick stencil via static-reduction. \\
\bottomrule
\end{tabular}
\caption{Hypotheses of \cref{thm:T9-informal} (Klein--Gordon residue
bound). The hypothesis footprint is identical in shape to T2's, which
strengthens the unified narrative ``Schr\"odinger $\subset$
Klein--Gordon on the same substrate with the same hypothesis
footprint''.}
\label{tab:T9-hyps}
\end{table}

\subsection{Ties to the non-relativistic theorems}

Two companion theorems exhibit the compatibility between T9 and the
earlier substrate facts:

\begin{itemize}[leftmargin=1.4em]
\item \code{kleinGordon\_mass\_shell\_consistent} ties the KG equation
to the relativistic mass shell $E^{2} = (pc)^{2} + (mc^{2})^{2}$ of
\code{DispersionFromLattice.lean}. The content of the lemma is the
identity $E^{2} = p^{2}c^{2} + m^{2}c^{4}$ (restated as
\code{relativisticEnergy\_sq\_eq}), which is the dispersion relation
a plane-wave solution of the continuum Klein--Gordon equation must
satisfy. The lattice Klein--Gordon residue vanishes as $\lP \to 0$
precisely when the plane-wave momentum $p$ and energy $E$ sit on this
shell.

\item \code{kleinGordon\_nonRelativistic\_limit} composes with
\cref{thm:T4-informal} to produce the analogous non-relativistic
limit for the Klein--Gordon frequency: subtracting the rest-energy
term $mc^{2}$ from the KG dispersion and approximating to order
$p^{4}/(4 m^{3} c^{2})$ recovers the Schr\"odinger kinetic term
$p^{2}/(2m)$. The ``Schr\"odinger $\subset$ Klein--Gordon $\subset$
relativistic'' hierarchy is therefore a chain of lattice-level
theorems.
\end{itemize}

The flat-vacuum specialisation
\code{coarseGrain\_satisfies\_kleinGordon\_dynamic\_flat} discharges
the \code{hcenter}/\code{hstencil} 9-point Gibbs hypotheses
automatically (by \code{informationDensityKL\_flat\_self}), so on
Minkowski the bound holds unconditionally for every mass and every
$(p, n)$. The static-regime reference form is
\code{coarseGrain\_satisfies\_kleinGordon\_static}.

\subsection{What Theorem~9 does and does not say}

Theorem~9 \emph{is} a machine-checked relativistic-to-discrete bridge:
the same discrete-gravity substrate that yields non-relativistic QM
as a chain of theorems (T1--T8) also yields the relativistic
Klein--Gordon equation, with an explicit Planck-scale UV cutoff as the
residue constant. The bridge extends the paper's claim from
``non-relativistic QM emerges from Omega-theory'' to
``non-relativistic QM and the relativistic Klein--Gordon equation
emerge from the same Omega-theory substrate, with their respective
residue bounds set by the Planck scale''.

Theorem~9 is \emph{not} a continuum Klein--Gordon equation. The
residue is bounded, not zero. Recovering zero residue requires the
$C^{4}$-bounded coarse-graining workstream. Theorem~9 is \emph{not}
QED or spinor QFT; the Dirac extension is open and stubbed in
\cref{sec:open} as item~\ref{open:dirac}. Theorem~9 is \emph{not} a
derivation of second quantisation, creation/annihilation operators, or
Fock space; the formalism is first-quantised lattice QM.

\section{Capstone --- Grand QM Emergence}\label{sec:capstone}

\subsection{The umbrella claim}\label{sec:capstone-claim}

\cref{sec:coarsegrain,sec:schrodinger,sec:born,sec:nonrel,sec:interference,sec:heisenberg,sec:measurement,sec:entanglement,sec:kleingordon}
deliver eight individual theorems (T1--T9; the Bell bundle T8a--T8d is
one theorem-slot; the Klein--Gordon relativistic bridge is T9), each a
machine-checked derivation of a standard-QM postulate or relativistic
residue bound from the V2 substrate. This section shows that these
are not nine independent results but a single composite statement:
for every dynamical snapshot sequence in the $\HasZF$ regime, for
every coarse-grained lattice field, for every finite region and every
tick, the \textbf{nine-conjunct postulate record}
\[
\text{states}+\text{superposition}+\text{Schr\"odinger}+\text{Born}+\text{interference}+\text{Heisenberg}+\text{measurement}+\text{entanglement}+\text{Klein--Gordon}
\]
holds simultaneously.

\subsection{The
\texorpdfstring{\code{QuantumMechanicsPostulates}}{QuantumMechanicsPostulates}
predicate}\label{sec:capstone-predicate}

\code{QuantumMechanicsCapstone.lean} introduces a structure
\begin{lstlisting}
structure QuantumMechanicsPostulates
    (d : DynamicalSnapshotSequence)
    (_hscope : d.HasZeroFunctional)
    (ψ : LatticeComplexField)
    (region : Finset LatticePoint)
    (tick : ℕ) : Prop where
  states         : ...                    -- P1: bounded ℓ² on region
  superposition  : ...                    -- P2: |ψ₁+ψ₂|² = ...
  schrodinger    : ...                    -- P3: Schrödinger bound
  born_rule      : ...                    -- P4: conservation
  interference   : ...                    -- P5: two-slit formula
  heisenberg     : ...                    -- P6: Δx·Δp ≥ ℏ/2
  measurement    : ...                    -- P7: 4-clause collapse
  entanglement   : IsEntangled bellField ∧
                     chshTsirelsonBell ... > 2   -- P8: Bell + CHSH
  kleinGordon    : ...                    -- P9: Klein--Gordon residue ≤ 16/ℓ_P² + m²c²/ℏ²
\end{lstlisting}

Each field is the $\Prop$-level statement corresponding to one of
the paper's theorems (P8 bundles the two Phase-6C headlines into a
single conjunct; P9 is the Klein--Gordon relativistic bridge of
\cref{sec:kleingordon}). P1--P6 depend on the specific state $\psi$,
region, and tick; P7--P8 quantify over outcome / angle as needed; P9
quantifies over rest mass $m : \Rbb$ and holds for every $m$,
including $m = 0$.

\subsection{The capstone theorem}\label{sec:capstone-theorem}

\begin{lstlisting}
/-- **THE GRAND QM EMERGENCE THEOREM.** Every postulate of
    non-relativistic quantum mechanics, plus the entanglement /
    CHSH postulate, is a machine-checked theorem of OmegaTheory V2
    in the static-functional regime. -/
theorem paper_grand_qm_emergence
    (d : DynamicalSnapshotSequence) (hscope : d.HasZeroFunctional)
    (ψ : LatticeComplexField) (region : Finset LatticePoint)
    (tick : ℕ) :
    QuantumMechanicsPostulates d hscope ψ region tick :=
  grand_qm_emergence d hscope ψ region tick
\end{lstlisting}

\paragraph{Provenance of each conjunct.}
The proof is a \emph{direct field-by-field citation} of the sibling
theorems---it introduces no new mathematical content.

\begin{longtable}{@{}p{0.22\linewidth} p{0.44\linewidth} p{0.28\linewidth}@{}}
\toprule
Conjunct & Source theorem & Source file \\
\midrule
\endhead
\code{states} & \code{coarseGrain\_info\_bounded} & \code{CoarseGrainingMap.lean} \\
\code{superposition} & \code{superposedField\_abs\_sq} & \code{Interference.lean} \\
\code{schrodinger} & \code{coarseGrain\_satisfies\_schrodinger\_dynamic} & \code{SchrodingerFromLattice.lean} \\
\code{born\_rule} & \code{bornRuleConservation} & \code{BornRule.lean} \\
\code{interference} & \code{two\_slit\_interference} & \code{Interference.lean} \\
\code{heisenberg} & \code{heisenberg\_uncertainty\_from\_lattice} & \code{Heisenberg.lean} \\
\code{measurement} & \code{measurement\_postulate} & \code{Measurement.lean} \\
\code{entanglement} & \code{bell\_inequality\_entanglement\_consistency} & \code{Entanglement.lean} \\
\code{kleinGordon} & \code{coarseGrain\_satisfies\_kleinGordon\_dynamic} & \code{KleinGordon.lean} \\
\bottomrule
\end{longtable}

\noindent No postulate is proved twice; the capstone's novelty is
organisational, not mathematical. The value is that \emph{one}
top-level statement now stands for the paper's claim, and downstream
consumers can cite a single theorem rather than seven.

\subsection{Minkowski specialisation}\label{sec:capstone-minkowski}

\begin{lstlisting}
theorem paper_grand_qm_emergence_on_minkowski
    (ψ : LatticeComplexField) (region : Finset LatticePoint) (tick : ℕ) :
    QuantumMechanicsPostulates minkowskiDynamicalSequence
      minkowskiDynamicalSequence_hasZeroFunctional ψ region tick :=
  grand_qm_emergence_on_minkowski ψ region tick
\end{lstlisting}

On the vacuum, $\HasZF$ is discharged automatically, so all nine
postulates hold for every $\psi$, every region, every tick, and every
rest mass $m : \Rbb$.

\subsection{What the capstone does and does not
say}\label{sec:capstone-neg}

It \emph{is} the single-statement formalisation of ``non-relativistic
quantum mechanics and the relativistic Klein--Gordon equation emerge
from Omega-Theory V2 on the same substrate'':
$\mathsf{QuantumMechanicsPostulates}$ is inhabited for every
dynamical-substrate / state / region / tick combination in the
$\HasZF$ regime, with the Klein--Gordon field P9 holding for every
rest mass $m : \Rbb$ (including $m = 0$). It is \emph{not} a new
physical mechanism: it introduces no mathematics beyond the nine
component theorems. It is \emph{not} the entirety of quantum
mechanics; extensions---$N$-body entanglement beyond two-body,
relativistic QFT with second quantisation, decoherence /
pointer-state reconstruction, spin-$1/2$ Dirac fermions, unified
QM-gravity---are out of scope and flagged in Section~\ref{sec:open}.

To the authors' knowledge, this is the first complete machine-checked
derivation of the postulate set of non-relativistic quantum mechanics
\emph{plus} the relativistic Klein--Gordon residue bound from a
discrete gravitational substrate in a proof assistant (best-effort
framing; no exhaustive-survey claim).

\section{Open Questions}\label{sec:open}

\begin{quote}
\textbf{Scope note.} The paper's scope is the nine-theorem chain
(T1--T9 of \cref{sec:coarsegrain,sec:schrodinger,sec:born,sec:nonrel,sec:interference,sec:heisenberg,sec:measurement,sec:entanglement,sec:kleingordon})
plus the capstone. Tracking state at the draft timestamp below:

\begin{itemize}
\item \textbf{Phase~1} (dynamical snapshot-sequence update rule)
  --- \emph{landed}. \code{SnapshotDynamics.lean}, by Polaris.
  Integrated in Section~\ref{sec:dynamical}.
\item \textbf{Phase~2} (dynamical Schr\"odinger bound) ---
  \emph{landed} under $\HasZF$.
  \code{coarseGrain\_satisfies\_schrodinger\_dynamic}, by
  \code{schrodinger\_prover}. Integrated as Section~\ref{sec:T2}.
\item \textbf{Phase~3} (Born rule as a conservation theorem) ---
  \emph{landed}. \code{BornRule.lean}, by
  \code{probability\_conservator}. Integrated as
  Section~\ref{sec:born}.
\item \textbf{Phase~4} (two-slit interference) --- \emph{landed}.
  \code{Interference.lean}, by \code{interference\_prover}.
  Integrated as Section~\ref{sec:interference}.
\item \textbf{Phase~6A} (Heisenberg uncertainty principle) ---
  \emph{landed}. \code{Heisenberg.lean}, by
  \code{uncertainty\_prover}. Integrated as
  Section~\ref{sec:heisenberg}.
\item \textbf{Phase~6B} (measurement / collapse postulate) ---
  \emph{landed}. \code{Measurement.lean}, by
  \code{measurement\_prover}. Integrated as
  Section~\ref{sec:measurement}.
\item \textbf{Phase~6C} (entanglement and Bell-inequality violation)
  --- \emph{landed}. \code{Entanglement.lean}, by
  \code{entanglement\_architect}. Integrated as
  Section~\ref{sec:entanglement}.
\item \textbf{Phase~7} (relativistic Klein--Gordon residue bound)
  --- \emph{landed}. \code{KleinGordon.lean}, by
  \code{klein\_gordon\_prover}. Integrated as
  Section~\ref{sec:kleingordon}; the Dirac extension is stubbed in
  \code{DiracOptional.lean} and catalogued as
  item~\ref{open:dirac} below.
\end{itemize}
\end{quote}

\begin{enumerate}
\item \textbf{Einstein side is not yet fully axiom-free.} The
  Einstein tensor emergence theorem
  \code{einstein\_with\_matter\_emergence} in
  \code{Emergence/EinsteinEmergence.lean} currently uses the HPW axiom
  as one step. HPW has been proved eliminable on an expanding family
  of regimes via the \code{HpwEliminableRegime} typeclass
  (\code{NOTES\_HPW\_ELIMINATION.md}): flat Minkowski, linearised
  metric, static-spherical vacuum, and (recently) a structural
  perturbed-flat path
  \code{hpw\_eliminable\_on\_perturbed\_flat} in
  \code{Emergence/HpwLinearised.lean}, which takes a
  \code{perturbedFlatInterpolant} from
  \code{Emergence/SmoothInterpolant.lean}---a $C^{4}$ interpolant for
  $g = \eta + \varepsilon h$ with $C^{4}$ bound $|\varepsilon|\cdot M_{h}$---
  and routes it through \code{linearisedDataOfPerturbedFlat} with
  minimal physical input (Laplacian and Ricci bounds plus a length
  scale). A companion construction \code{slowlyVaryingInterpolant}
  shows that every smooth $g_{\mathrm{cont}}$ with a uniform $C^{4}$
  bound is its own interpolant on the lattice it samples; the
  curried adapter \code{SmoothInterpolantData.toSmoothMetricField}
  auto-fills the $g_{\mathrm{cont}}$ / $h_{\mathrm{interpolates}}$
  fields for any physics model that builds its discrete metric via
  \code{DiscreteMetric.ofContinuum}. What remains open is a
  \emph{generic} existence theorem (e.g.\ via a bump-convolution
  Whitney-extension argument) on arbitrary smooth backgrounds; this
  is explicitly deferred in the header of
  \code{SmoothInterpolant.lean} and currently blocks on Whitney
  extension infrastructure not yet present in Mathlib. The QM-bridge
  theorems of the present paper do not depend on HPW, so this is a
  limitation on the \emph{gravitational} side of the same substrate,
  not on the wave-mechanical side we derive here.

\item \textbf{Geometrically-sourced phase is not derived.} The
  phase-carrying map \code{coarseGrainWithPhase} and its plane-wave
  adapter accept a phase function as input. Phase~4 (interference)
  specialises to a concrete plane-wave phase, which is enough for the
  superposition / two-slit theorem but does not constitute a
  derivation of the phase from geometry. A principled derivation---
  e.g.\ from the \code{connectionForm} or \code{spinTorsion} of the
  lattice geometry---would close the Aharonov--Bohm case and several
  other gauge-sensitive predictions.

\item \textbf{Unconditional dynamical Schr\"odinger bound.} Phase~2
  closes the dynamical Schr\"odinger-shape inequality under $\HasZF$.
  Off-regime, closing the bound requires a \textbf{KL-linearisation
  bridge}
  $\KL(g+\delta g)\approx\KL(g)+\langle\partial\KL/\partial g,\delta g\rangle$
  not yet formalised in V2. The Phase-2 brief forbids axioms; stating
  an unconditional dynamical bound without the linearisation bridge
  would therefore require introducing one. The honest current
  frontier is the $\HasZF$-scoped bound.

\item \textbf{Phase~1's $F=\Delta_{\mathrm{lat}}(\cdot)$ update is
  structurally motivated, not derived from the healing-flow PDE.}
  The dynamical update rule of Section~\ref{sec:dynamical} is
  dictated by the requirement that the coarse-grained evolution
  carry a Schr\"odinger $-i\hbar\,\Delta/2m$ coefficient on the QM
  side. It is not derived from a weak-field limit of
  \code{OmegaTheory.HealingFlow.Flow}. A follow-up workstream should
  prove that the Phase~1 update is the weak-field / short-time limit
  of the full healing-flow PDE. This is the single most important
  structural open question on the substrate side.

\item \textbf{Lattice-Fourier eigenvalue extraction is not yet
  formalised.} Section~\ref{sec:nonrel-bridge} states the plane-wave
  identities needed to chain Theorems~\ref{thm:T2-informal} and
  \ref{thm:T4-informal}; those are pure lattice-Fourier facts, each a
  few trigonometric lemmas. None is in the Lean development at the
  time of writing. Concretely, the missing theorems are
  (i) \code{discreteLaplacianC\_planeWave\_eigenvalue},
  (ii) \code{planeWave\_eigenvalue\_Taylor}, and
  (iii) a paper wrapper binding $K:=\lambda(k,\lP)$ into
  \code{paper\_schrodinger\_is\_nonrel\_limit}.

\item \textbf{The mass-shell derivation is tick-counting, not
  stress-energy.} A more principled derivation from the stress-energy
  tensor of the healing flow is plausible but not written.

\item \textbf{Empirical falsifiability is in principle, not in
  practice.} \code{NOTES\_QM\_AS\_DISCRETE\_GRAVITY.md} §5
  \cite{NotesQMasDiscreteGravity} documents this honestly: the
  predicted power-law gate-fidelity curve $F(T)=F_{0}/(1+\alpha T)$
  differs from the Arrhenius form $F_{0}\exp(-E/\kB T)$ by a
  functional-form signature, but the numerical coupling
  $\alpha=\kB\tP/(2\hbar)\approx 3.5\times 10^{-33}\,\mathrm{K}^{-1}$
  is 28 orders of magnitude below current cryogenic-gate precision.

\item \textbf{Klein--Gordon continuum limit (zero-residue
  equation).}\label{open:kg-coarse}
  \cref{sec:kleingordon} proves the discrete Klein--Gordon
  residue bound $\|\mathrm{KG}[\psi](p, n)\| \le 16/\lP^{2} +
  m^{2} c^{2}/\hbar^{2}$. The continuum relation
  $(\Box + m^{2} c^{2}/\hbar^{2})\psi = 0$ is the $\lP \to 0$
  limit of this, not the discrete equation itself. Recovering
  \emph{zero} residue (rather than a Planck-scale UV-cutoff bound)
  requires a $C^{4}$-bounded coarse-graining over many lattice
  sites together with a Taylor-bound lemma for the second time
  difference and the five-point Laplacian. This is pure
  lattice-Fourier / finite-differences analysis, no new physics; it
  remains the natural continuum-limit follow-up to Theorem~9.

\item \textbf{Spin-$1/2$ extension via Cartan--Sciama--Kibble
  torsion.}\label{open:dirac}
  \code{OmegaTheory/Emergence/DiracOptional.lean} introduces a
  Prop-carrying structure \code{DiracFromLatticeData} that records
  what a full lattice Dirac derivation must supply: a Clifford
  algebra $\mathrm{Cl}(1, 3)$ with concrete gamma matrices, a lattice
  discretisation compatible with chiral symmetry (Nielsen--Ninomiya
  evasion via Wilson fermions, staggered fermions, or domain-wall
  fermions), and a Cartan--Sciama--Kibble spin-torsion coupling via
  the existing \code{OmegaTheory.Torsion.SpinTorsion}. The
  \code{waveEquation} predicate is currently abstract; the
  \code{massShellConsistent} and \code{nonRelativisticLimit} fields
  are already theorems (inherited from \code{SpecialRelativity} and
  \code{KleinGordon}). A trivial inhabitant
  \code{diracFromLattice\_partial} is provided as a stub,
  acknowledging that a full lattice Dirac discretisation (proper
  gamma-matrix algebra, chiral-symmetry preservation, and
  spin-torsion coupling) remains an explicit follow-up workstream
  --- the spin-$1/2$ analogue of the scalar Klein--Gordon bridge of
  \cref{sec:kleingordon}.

\item \textbf{Many-body states beyond two-body Bell states.}
  Phase~6C delivers the two-body Bell state $\BellField$ and its
  CHSH-violation signature, but the present formalisation does not
  cover arbitrary N-body entangled states, reduced-density-matrix
  dynamics, or multi-body monogamy-of-entanglement statements.

\item \textbf{Continuous-time limit.} All theorems are discrete-tick
  statements. The smooth $\tP\to 0$ limit---recovering the continuous
  Schr\"odinger PDE rather than its discrete approximant---is left
  as future work.

\item \textbf{Second quantisation and QFT.} Theorem~9 delivers the
  relativistic scalar Klein--Gordon residue at the first-quantised
  lattice level. Creation and annihilation operators, Fock space,
  causal commutation relations, and propagator machinery are not
  derived. QFT-on-the-substrate is a natural sequel to a zero-residue
  continuum KG limit (item~\ref{open:kg-coarse}).
\end{enumerate}

None of these items invalidates Phases~1--6 or Theorem~9. They
delimit the claim. The paper's claim is: given the discrete-gravity
healing-flow substrate of Omega-Theory V2, \emph{all nine} of (T1--T7
as before:) Schr\"odinger evolution, Born-rule probability
conservation, relativistic-to-non-relativistic dispersion, two-slit
interference, the Heisenberg uncertainty principle, the measurement /
collapse postulate, and Tsirelson-bound-attaining entanglement with
CHSH-inequality violation, \emph{plus} (T9:) the relativistic
Klein--Gordon residue bound with an explicit Planck-scale UV cutoff
$16/\lP^{2} + m^{2}c^{2}/\hbar^{2}$, are \emph{theorems}, not
postulates. It is not: Omega-Theory V2 reproduces every phenomenon
of quantum mechanics, or of quantum field theory.

The value of distinguishing these two claims is methodological. A
proof-assistant-verified partial result is, in our view, more useful
to the community than an unverified complete one, because each open
item above is now a concrete next target rather than a vague research
programme---and the boundary of ``what is proved today'' is exactly
where a reviewer can focus their critical attention.

\section*{Acknowledgments}\label{sec:ack}
\addcontentsline{toc}{section}{Acknowledgments}

This formalization was produced by a team of AI agents operating on
the Omega-Theory V2 codebase, under the direction of the human author.
The individual contributions are:
\begin{itemize}
\item \textbf{Altair} --- coarse-graining map and phase-accepting
  extension in \code{CoarseGrainingMap.lean}, including the defining
  identity \code{sq\_abs\_coarseGrain} and the finite-region
  $\ell^{2}$ bound.
\item \textbf{Sirius} (Phase~2 lead) --- \code{SchrodingerFromLattice.lean}
  including the complex discrete Laplacian, the static-regime
  headline bound \code{coarseGrain\_satisfies\_schrodinger\_static},
  the phase-aware specialisation, and the Phase-2 dynamical lift
  \code{coarseGrain\_satisfies\_schrodinger\_dynamic} together with
  flat-background and phase-aware corollaries.
\item \textbf{Bridger} --- \code{DispersionBridge.lean} including the
  closed-form quartic bound \code{nonRelativistic\_energy\_approx}
  and the \code{schrodinger\_is\_nonrel\_limit} bridge lemma.
\item \textbf{Polaris} (Phase~1 lead) --- \code{SnapshotDynamics.lean}
  including the \code{DynamicalSnapshotSequence} structure, the
  Laplacian-of-metric update rule, the
  \code{metric\_update\_linear\_in\_t\_P} control, and the Minkowski
  vacuum-consistency theorem.
\item \textbf{\code{probability\_conservator}} (Phase~3 lead) ---
  \code{BornRule.lean} including the tick-invariance key lemma, the
  headline conservation theorem \code{bornRuleConservation}, the
  iterated form, the Minkowski specialisation, and the honest
  off-regime residue bound.
\item \textbf{\code{interference\_prover}} (Phase~4 lead) ---
  \code{Interference.lean} including the complex-field superposition,
  the plane-wave coarse-grained field, the algebraic heart
  \code{superposedField\_abs\_sq}, the headline
  \code{two\_slit\_interference} and its Boltzmann-weight form, the
  canonical visibility, the constructive peak / destructive null
  corollaries, and the Phase-3 compatibility corollary.
\item \textbf{Bellatrix} (Phase~6A lead, \code{uncertainty\_prover})
  --- \code{Heisenberg.lean} including the lattice position and
  momentum operators, the $L^{2}$-variance definition, the
  Robertson-cross-term machinery, the exact Robertson inequality, the
  $\CMM$ / \code{CommutatorMatchesValue} predicates, the abstract
  uncertainty theorem, and the final position-momentum theorem.
\item \textbf{\code{measurement\_prover}} (Phase~6B lead) ---
  \code{Measurement.lean} including the pointwise projector and its
  idempotence, the region-relative norm, the renormaliser, the
  composition \code{postMeasurementState}, the probability
  functional, the unit-norm corollary, the support-concentration
  lemma, the headline four-clause bundle, the Phase-3 bridge, the
  crucial non-unitarity witness
  \code{postMeasurement\_not\_in\_coarseGrain\_image}, and the
  Minkowski specialisation.
\item \textbf{\code{entanglement\_architect}} (Phase~6C lead) ---
  \code{Entanglement.lean} including the two-body field type, the
  tensor-product operation with bilinearity and abs-sq factorisation
  lemmas, the \code{IsProduct} / \code{IsEntangled} predicates, the
  explicit Bell-state construction, the structural entanglement
  theorem, the correlator and the headline cosine identity, the CHSH
  functional, the Tsirelson-bound attainment, the classical-Bell-bound
  violation, and the consistency bundle.
\item \textbf{Saiph} (Phase~6D capstone lead) ---
  \code{QuantumMechanicsCapstone.lean} including the
  \code{QuantumMechanicsPostulates} 9-field structure, the
  \code{grand\_qm\_emergence} umbrella theorem, the Minkowski
  specialisation, and the named projection corollaries; the
  integration of the P9 \code{kleinGordon} field after Theorem~9
  landed.
\item \textbf{\code{interpolant\_extender}} --- the
  perturbed-flat HPW-elimination path:
  \code{SmoothInterpolant.lean} (\code{perturbedFlatInterpolant} for
  $g = \eta + \varepsilon h$, \code{slowlyVaryingInterpolant} as the
  self-interpolant for any uniformly-$C^{4}$-bounded
  $g_{\mathrm{cont}}$, and the curried adapter
  \code{SmoothInterpolantData.toSmoothMetricField}) and
  \code{HpwLinearised.lean}
  (\code{linearisedDataOfPerturbedFlat}, the end-to-end axiom-free
  \code{hpw\_eliminable\_on\_perturbed\_flat}). This extends the HPW
  elimination story from the three earlier regimes to any physics
  model that builds its discrete metric via
  \code{DiscreteMetric.ofContinuum}.
\item \textbf{\code{klein\_gordon\_prover}} (Phase~7 lead) ---
  \code{KleinGordon.lean} including the discrete second-time
  difference, the Klein--Gordon residue definition, the
  \code{kleinGordonBoundConst} constant
  $16/\lP^{2} + m^{2}c^{2}/\hbar^{2}$, the
  \code{coarseGrain\_satisfies\_kleinGordon\_static} and dynamical
  bounds, the Minkowski flat-vacuum corollaries, and the
  companion lemmas \code{kleinGordon\_mass\_shell\_consistent} and
  \code{kleinGordon\_nonRelativistic\_limit} that tie T9 to T4's
  dispersion bridge and T2's Schr\"odinger kinetic term.
  Also \code{DiracOptional.lean} including the Prop-carrying
  \code{DiracFromLatticeData} structure and the
  \code{diracFromLattice\_partial} stub recording the open
  spin-$1/2$ extension via Cartan--Sciama--Kibble torsion.
\item \textbf{Alnilam} (\code{paper\_draft}) ---
  \code{QmBridgePaper.lean} (24 paper-level wrapper theorems across
  Theorems~1--8 plus the capstone) and the first version of the
  present manuscript.
\item \textbf{\code{paper\_polisher}} --- the LaTeX submission
  package (\texttt{main.tex}, \texttt{letter.tex},
  \texttt{cover\_letter.tex}, \texttt{refs.bib},
  \texttt{README\_SUBMISSION.md}), including the integration of
  Theorem~9 (\cref{sec:kleingordon}) and the Dirac stub
  (\cref{sec:open} item~\ref{open:dirac}) into the manuscript after
  \code{klein\_gordon\_prover} landed Phase~7.
\end{itemize}

Upstream infrastructure was built by earlier agents in the project
(\textbf{Vega}, \textbf{Rigel}, and others credited in
\code{README.md}). The underlying physical theory, the Lean project
scaffolding, and editorial direction are due to Norbert Marchewka. Any
errors are ours.

This work was carried out without external grant funding. The Lean
proofs have been verified against Mathlib v4.29.0. The source is
published with the paper.

\bibliographystyle{plain}
\bibliography{refs}

\end{document}
